\documentclass[10pt,a4paper]{article}
\usepackage[utf8]{inputenc}
\usepackage[margin=2.2cm]{geometry}
\usepackage{amsmath,amssymb,amsthm,mathtools}
\usepackage{microtype}
\usepackage{booktabs}
\usepackage{enumitem}
\usepackage{xcolor}
\usepackage[colorlinks=true,linkcolor=blue!50!black,citecolor=blue!50!black]{hyperref}
\usepackage[spanish,ruled,vlined,linesnumbered,noend]{algorithm2e}

% ---------------------------------------------------------------- algorithm2e
\renewcommand{\algorithmcfname}{Algoritmo}
\SetKwInput{KwIn}{Entrada}
\SetKwInput{KwOut}{Salida}
\SetKwProg{Fn}{función}{}{}
\SetKw{KwRet}{devolver}
\SetKw{KwTo}{hasta}
\SetKw{KwBreak}{cortar}
\SetKw{KwCont}{continuar}
\SetKw{KwAnd}{y}
\SetKw{KwOr}{o}
\SetKw{KwNot}{no}
\SetKwFor{ForEach}{para cada}{hacer}{}
\SetKwFor{For}{para}{hacer}{}
\SetKwFor{While}{mientras}{hacer}{}
\SetKwIF{If}{ElseIf}{Else}{si}{entonces}{si no, si}{si no}{}
\SetCommentSty{textit}
\SetKwComment{tcp}{$\triangleright$\ }{}
\SetAlgoNlRelativeSize{-1}
\setlength{\algomargin}{1.6em}

\SetKwFunction{SotSort}{SotSort}
\SetKwFunction{Ordenado}{YaOrdenado}
\SetKwFunction{PocosGiros}{PocosGiros}
\SetKwFunction{Natural}{MergeNatural}
\SetKwFunction{Potencia}{Potencia}
\SetKwFunction{FusTope}{FusionarTope}
\SetKwFunction{Respaldo}{MergeRespaldo}
\SetKwFunction{Quick}{QuickEstable}
\SetKwFunction{Pivote}{Pivote}
\SetKwFunction{Part}{ParticiónEstable}
\SetKwFunction{PartBloque}{ParticiónBloque}
\SetKwFunction{Fusionar}{Fusionar}
\SetKwFunction{Rotar}{Rotar}
\SetKwFunction{TRF}{TripleReversiónFusionada}
\SetKwFunction{Hoja}{Hoja}
\SetKwFunction{PingPong}{PingPongPot2}
\SetKwFunction{Parity}{ParityMerge}
\SetKwFunction{MIzq}{MergeIzq}
\SetKwFunction{MDer}{MergeDer}
\SetKwFunction{SymMerge}{SymMerge}
\SetKwFunction{Insercion}{Inserción}
\SetKwFunction{FinAsc}{FinAsc}
\SetKwFunction{FinDesc}{FinDesc}
\SetKwFunction{Aleatorio}{Aleatorio}
\SetKwFunction{CoRango}{CoRango}
\SetKwFunction{Conteo}{Conteo}
\SetKwFunction{Disp}{Disponible}
\SetKwFunction{Galope}{Galope}
\SetKwFunction{Paralelo}{Paralelo}
\SetKwFunction{PlanF}{Plan}
\SetKwFunction{PorClave}{OrdenPorClave}
\SetKwFunction{Rapido}{OrdenRápido}
\SetKwFunction{Reservar}{Reservar}
\SetKwFunction{PartVec}{PartBloqueVec}

% ---------------------------------------------------------------- teoremas
\theoremstyle{plain}
\newtheorem{teorema}{Teorema}[section]
\newtheorem{lema}[teorema]{Lema}
\newtheorem{corolario}[teorema]{Corolario}
\theoremstyle{definition}
\newtheorem{definicion}[teorema]{Definición}
\newtheorem{supuesto}{Supuesto}
\theoremstyle{remark}
\newtheorem{observacion}[teorema]{Observación}
\renewcommand{\proofname}{Demostración}

% ---------------------------------------------------------------- notación
\newcommand{\lt}{\prec}
\newcommand{\nlt}{\not\prec}
\newcommand{\les}{\preccurlyeq}
\newcommand{\eqv}{\sim}
\newcommand{\lts}{\prec_{s}}
\newcommand{\cat}{\mathbin{+\!\!+}}
\newcommand{\buf}{\mathit{buf}}
\newcommand{\Bsz}{B}
\newcommand{\leaf}{L}
\newcommand{\etq}{\tau}
\newcommand{\elems}{\operatorname{elem}}
\newcommand{\filt}{\operatorname{filtro}}
\newcommand{\Mrg}{\mathcal{M}}
\newcommand{\rev}{\operatorname{rev}}
\newcommand{\hoare}[3]{\{\,#1\,\}\;\;#2\;\;\{\,#3\,\}}
\newcommand{\seg}[3]{\ensuremath{#1[#2\mathinner{.\,.}#3)}}
\newcommand{\sortd}{\operatorname{ord}}
\newcommand{\coh}{\operatorname{coh}}

\title{\textbf{SotSort}\\[3pt]\large Pseudocódigo, demostraciones de corrección (lógica de Hoare)\\ y análisis de complejidad --- versiones 6.1 a 8-r}
\author{Lucas Sotomayor}
\date{17 de noviembre de 2026 --- borrador para revisión}

\begin{document}
\maketitle

\begin{abstract}
SotSort es un algoritmo de ordenamiento estable, sin recursión y con memoria
auxiliar constante: un buffer de $\Bsz=512$ elementos y pilas de tamaño acotado.
Combina un quicksort estable cuya partición se construye por bloques y rotaciones,
hojas ordenadas por \emph{parity merges} y, para datos con estructura, un merge por
runs naturales con el criterio de powersort. Este documento da el pseudocódigo
completo de la versión 6.1 y demuestra, para cada procedimiento, su corrección
total como tripla de Hoare (invariantes de lazo y variantes para la terminación), la
estabilidad del resultado, la cota de memoria, la profundidad de cada pila y cotas
de peor caso para movimientos y comparaciones.
\end{abstract}

\tableofcontents

%======================================================================
\section{Alcance, supuestos y una corrección surgida de la prueba}

\paragraph{Qué se demuestra.} Para el pseudocódigo de este documento:
(i) corrección total de cada procedimiento como tripla de Hoare, con invariantes y
variantes; (ii) que \SotSort devuelve la \emph{única} permutación estable ordenada de
la entrada (Teorema~\ref{thm:main}); (iii) memoria auxiliar constante y cotas de cada
pila; (iv) cotas de peor caso: $O(n\log^2 n)$ movimientos y
$O\!\big(n\log n + \tfrac{n}{\Bsz}\log^3 n\big)$ comparaciones.

\paragraph{Qué \emph{no} se demuestra.}
\begin{itemize}[nosep]
\item La equivalencia entre este pseudocódigo y la implementación en C++
  (\texttt{stable\_qsort.hpp}). Esa correspondencia está respaldada por pruebas
  (fuzzing contra \texttt{std::stable\_sort}: 6557 casos sin fallas con
  gcc+ASan/UBSan y con clang), no por una demostración.
\item Ninguna afirmación de rendimiento real. Las mediciones contra blitsort son
  experimentales.
\item Las demostraciones son manuales y no están verificadas mecánicamente. Un paso
  natural sería formalizarlas en Why3, Dafny o Coq.
\end{itemize}

\paragraph{Nota de esta edición (v8-r).} Las secciones 1 a 11 describen la versión 6.1
con $\Bsz=512$. La sección~\ref{sec:rapido} agrega el modo rápido de la versión 8. La
sección~\ref{sec:v8r} agrega los cambios de la versión 7 (buffer propio de unos 16\,KB,
pivote con posiciones pseudoaleatorias, partición vectorial) y los de la reconstrucción
v8-r; la sección~\ref{sec:ooo}, un planificador experimental. El archivo
\texttt{stable\_qsort.hpp} y las cifras de pruebas que se mencionan abajo son de las
versiones originales, cuyo código se perdió; las cifras vigentes están en la
sección~\ref{sec:v8r-pruebas}.

\begin{supuesto}\label{sup:orden}
$\lt$ es un orden débil estricto (Def.~\ref{def:swo}) sobre el tipo de los elementos,
y su evaluación es total, determinista y sin efectos laterales.
\end{supuesto}
\begin{supuesto}\label{sup:ident}
Mover un elemento preserva su identidad. Las \emph{copias} (pivote $p$, cota inferior
$\mathit{lb}$, muestras) solo se usan como argumento de $\lt$, nunca se escriben en
el arreglo.
\end{supuesto}
\begin{supuesto}\label{sup:n}
$n < 2^{63}$, constantes $\Bsz = 512$, $\leaf = 128 \le \Bsz$ y $N_9 = 8192$.
\end{supuesto}

\paragraph{Corrección surgida de la demostración (v6 $\to$ v6.1).}
Al demostrar la cota de profundidad de \Quick (Lema~\ref{lem:prof}) apareció un hueco.
El paso que separa el bloque de elementos equivalentes al pivote no consumía
presupuesto del respaldo introsort. Un adversario podía encadenar muchos de esos
pasos, cada uno quitando pocos elementos, y la profundidad dejaba de ser $O(\log n)$.
En v6.1 ese paso descuenta presupuesto cuando el bloque separado tiene menos de
$1/16$ del rango (líneas 12 y 18 del Algoritmo~\ref{alg:quick}). El cambio se aplicó
también al código C++: los tests siguen sin fallas y el rendimiento quedó dentro del
ruido de medición.

%======================================================================
\section{Modelo y lemas base}

\begin{definicion}[Orden débil estricto]\label{def:swo}
Una relación $\lt$ es un \emph{orden débil estricto} si es irreflexiva
($x\nlt x$), transitiva y la \emph{incomparabilidad}
$x\eqv y \iff x\nlt y \wedge y\nlt x$ es transitiva. Se escribe
$x\les y \iff y\nlt x$.
\end{definicion}

\begin{lema}[Propiedades]\label{lem:swo}
Si $\lt$ es un orden débil estricto, entonces:
\begin{enumerate}[label=(\alph*),nosep]
\item $\eqv$ es una relación de equivalencia;
\item (transitividad negativa) $x\lt z \Rightarrow x\lt y \vee y\lt z$ para todo $y$;
\item $x\lt y \wedge y\les z \Rightarrow x\lt z$ \quad y \quad $x\les y\wedge y\lt z \Rightarrow x\lt z$;
\item $\les$ es transitiva y total;
\item $x\eqv y \Rightarrow (x\lt p \Leftrightarrow y\lt p)\wedge(p\lt x\Leftrightarrow p\lt y)$.
\end{enumerate}
\end{lema}
\begin{proof}
(a) Reflexiva por irreflexividad, simétrica por definición y transitiva por hipótesis.
(b) Supongamos $x\lt z$, $x\nlt y$ y $y\nlt z$. Si fuera $y\lt x$, por transitividad
$y\lt z$, contradicción; luego $x\eqv y$. Si fuera $z\lt y$, entonces
$x\lt z\lt y$ da $x\lt y$, contradicción; luego $y\eqv z$. Por (a), $x\eqv z$, lo que
contradice $x\lt z$.
(c) Para la primera, aplicando (b) a $x\lt y$ con intermedio $z$ se obtiene
$x\lt z\vee z\lt y$; como $z\nlt y$, queda $x\lt z$. La segunda es simétrica.
(d) Transitividad: si $y\nlt x$, $z\nlt y$ y $z\lt x$, (b) da $z\lt y\vee y\lt x$,
contradicción. Totalidad: si $x\not\les y$ entonces $y\lt x$, y por asimetría
($x\lt y\wedge y\lt x\Rightarrow x\lt x$) $x\nlt y$, es decir, $y\les x$.
(e) Si $x\lt p$ y $y\nlt p$, entonces $p\les y$ y por (c) $x\lt y$, lo que contradice
$x\eqv y$. El resto es análogo.
\end{proof}

\begin{definicion}[Etiquetas, orden estable y salida esperada]\label{def:omega}
Sea $A_0[0..n)$ la entrada. El elemento $A_0[i]$ lleva la \emph{etiqueta}
$\etq(A_0[i]) = i$, que lo acompaña en cada movimiento (Supuesto~\ref{sup:ident}).
Se define
\[
e \lts f \iff e\lt f \;\vee\; (e\eqv f \wedge \etq(e)<\etq(f)).
\]
Por el Lema~\ref{lem:swo}(d), $\lts$ es un orden total estricto sobre los elementos
de $A_0$. Se denota $\Omega$ a la única disposición de esos elementos creciente según
$\lts$. La especificación de un ordenamiento estable es: \emph{la salida es
$\Omega$}.
\end{definicion}

\begin{definicion}[Ordenada, coherente]
Una secuencia $S=\langle s_0,\dots,s_{k-1}\rangle$ está \emph{ordenada}, y se escribe
$\sortd(S)$, si $\forall i<j:\ s_j\nlt s_i$. Es \emph{coherente}, $\coh(S)$, si
$\forall i<j:\ s_i\eqv s_j\Rightarrow\etq(s_i)<\etq(s_j)$. Se escribe $\elems(S)$
para el conjunto de sus elementos; es un conjunto porque las etiquetas son distintas.
\end{definicion}

\begin{lema}\label{lem:coh}
$S$ es creciente según $\lts$ si y solo si $\sortd(S)\wedge\coh(S)$. En particular,
si $\elems(S)=\elems(\Omega[I])$ para un intervalo $I$ y vale
$\sortd(S)\wedge\coh(S)$, entonces $S=\Omega[I]$.
\end{lema}
\begin{proof}
($\Rightarrow$) Es inmediato. ($\Leftarrow$) Para $i<j$, si $s_i\lt s_j$ no hay nada
que probar. Si no, como $s_j\nlt s_i$, se tiene $s_i\eqv s_j$ y la coherencia da
$\etq(s_i)<\etq(s_j)$. Para la segunda parte: $\Omega[I]$ es creciente según $\lts$,
y un conjunto finito tiene una única disposición creciente bajo un orden total.
\end{proof}

\begin{definicion}[Operación OE]
Una transformación $S\mapsto S'$ con $\elems(S')=\elems(S)$ \emph{preserva el orden
de equivalentes} (es OE) si, para todo par $x\eqv y$, que $x$ preceda a $y$ en $S$
implica que $x$ preceda a $y$ en $S'$.
\end{definicion}

\begin{lema}\label{lem:oe}
(a) Una operación OE preserva la coherencia. (b) La composición de operaciones OE es
OE. (c) Una operación que reordena un subrango $[l,r)$ y deja fijo el resto es OE
sobre el arreglo si y solo si lo es sobre $[l,r)$.
\end{lema}
\begin{proof}
(a) y (b) son inmediatas de la definición. (c) El orden relativo entre un elemento
de adentro y uno de afuera no cambia, porque las posiciones de afuera quedan fijas y
las de adentro siguen siendo un mismo intervalo.
\end{proof}

\begin{lema}[Rotación segura]\label{lem:rotsegura}
Pasar de $X\cat Y$ a $Y\cat X$ es OE si y solo si no existen $x\in X$, $y\in Y$ con
$x\eqv y$.
\end{lema}
\begin{proof}
Los pares que están ambos en $X$ o ambos en $Y$ conservan su orden. Los únicos pares
que se invierten son los de la forma $(x,y)$ con $x\in X$, $y\in Y$.
\end{proof}

\begin{definicion}[Filtro y partición estable]
$\filt(P,S)$ es la subsecuencia de los elementos de $S$ que cumplen $P$. La partición
estable es $\mathrm{Part}_P(S)=\filt(P,S)\cat\filt(\neg P,S)$. Un predicado $P$ es
\emph{$\eqv$-invariante} si $x\eqv y\Rightarrow P(x)=P(y)$.
\end{definicion}

\begin{lema}\label{lem:part}
(a) Si $P$ es $\eqv$-invariante, $S\mapsto\mathrm{Part}_P(S)$ es OE.
(b) Los predicados $P_<(x)\equiv x\lt p$ y $P_\le(x)\equiv p\nlt x$ son
$\eqv$-invariantes.
(c) Si $\sortd(S)$, entonces $S=\mathrm{Part}_{P_<}(S)=\mathrm{Part}_{P_\le}(S)$.
\end{lema}
\begin{proof}
(a) Dos elementos equivalentes terminan del mismo lado y cada lado conserva el orden
de $S$. (b) Es el Lema~\ref{lem:swo}(e). (c) Sean $i<j$ con $P_<(s_j)$ y
$\neg P_<(s_i)$, es decir, $s_j\lt p$ y $p\les s_i$. Por el Lema~\ref{lem:swo}(c),
$s_j\lt s_i$, lo que contradice $\sortd(S)$. Por lo tanto los elementos que cumplen
$P_<$ preceden a los que no. Para $P_\le$: con $p\nlt s_j$ y $p\lt s_i$, el
Lema~\ref{lem:swo}(c) da $s_j\lt s_i$, la misma contradicción.
\end{proof}

\begin{definicion}[Merge estable]\label{def:merge}
Dadas $X$ e $Y$ ordenadas, $\Mrg(X,Y)$ es el único entrelazado $Z$ de $X$ e $Y$ que
preserva el orden interno de cada una, está ordenado y en el que, si $x\in X$,
$y\in Y$ y $x\eqv y$, $x$ precede a $y$.
\end{definicion}

\begin{lema}[Merge]\label{lem:merge}
Sean $X[0..a)$ e $Y[0..b)$ ordenadas y $Z=\Mrg(X,Y)$.
\begin{enumerate}[label=(\alph*),nosep]
\item \emph{Paso de cabeza.} Si $Z[0..k)=\Mrg(X[0..h_x),Y[0..h_y))$ con $h_x<a$ y
  $h_y<b$, entonces $Z[k]=Y[h_y]$ si $Y[h_y]\lt X[h_x]$ y $Z[k]=X[h_x]$ en caso
  contrario.
\item \emph{Paso de cola.} Si $Z[N-k..N)$, con $N=a+b$, es el merge de los últimos
  $c_x$ elementos de $X$ y los últimos $c_y$ de $Y$, con $t_x=a-1-c_x\ge 0$ y
  $t_y=b-1-c_y\ge 0$, entonces $Z[N-1-k]=X[t_x]$ si $Y[t_y]\lt X[t_x]$ y
  $Z[N-1-k]=Y[t_y]$ en caso contrario.
\item \emph{Existencia y unicidad.} $\Mrg(X,Y)$ existe y es único. Además, cada
  segmento $Z[i..j)$ es $\Mrg$ de los segmentos de $X$ e $Y$ que contiene.
\item \emph{Estabilidad.} Si $X\cat Y$ está ordenada por partes (es decir, $X$ e $Y$
  están ordenadas), reemplazar $X\cat Y$ por $\Mrg(X,Y)$ es OE.
\end{enumerate}
\end{lema}
\begin{proof}
(a) $Z[k]$ es el primer elemento no consumido. Los candidatos son las cabezas $X[h_x]$
e $Y[h_y]$, porque cada secuencia está ordenada y su orden se preserva. Si
$Y[h_y]\lt X[h_x]$, poner primero $X[h_x]$ violaría el orden; si $X[h_x]\lt Y[h_y]$,
lo inverso; si $X[h_x]\eqv Y[h_y]$, la regla de desempate exige $X[h_x]$ primero.
(b) Es simétrico: el último no consumido es $X[t_x]$ si es estrictamente mayor y, ante
equivalencia, $Y[t_y]$ va después.
(c) La existencia se obtiene por inducción aplicando (a). La unicidad, porque (a)
determina cada posición. Un segmento de un entrelazado ordenado con la regla de
desempate cumple las mismas propiedades.
(d) Los pares equivalentes dentro de $X$ o de $Y$ conservan su orden. Un par $x\in X$,
$y\in Y$ con $x\eqv y$ tenía a $x$ antes, porque $X$ está a la izquierda, y $\Mrg$ lo
mantiene antes.
\end{proof}

\paragraph{Triplas de Hoare.} $\hoare{P}{S}{Q}$ indica \emph{corrección total}: si
$P$ vale antes de ejecutar $S$, entonces $S$ termina y $Q$ vale al terminar. $A_1$
denota el contenido del arreglo al entrar al procedimiento. ``$A\equiv A_1$ fuera de
$[l,r)$'' es la condición de marco: el procedimiento solo escribe en $[l,r)$ y en
$\buf$. Para cada lazo se da un invariante $I$ y una variante $V\in\mathbb{N}$ que
decrece estrictamente en cada iteración.

%======================================================================
\section{Rotación}

\begin{algorithm}[H]
\caption{\textsc{Rotar} y \textsc{TripleReversiónFusionada}}\label{alg:rotar}
\Fn{\Rotar{$l, m, r$}}{
  $a \gets m - l$;\quad $b \gets r - m$\;
  \lIf{$a = 0$ \KwOr $b = 0$}{\KwRet}
  \lIf{$a \le \Bsz$}{$\buf[0..a)\gets A[l..m)$; mover $A[m..r)\to A[l..l{+}b)$ (de izquierda a derecha); $A[l{+}b..r)\gets\buf[0..a)$; \KwRet}
  \lIf{$b \le \Bsz$}{$\buf[0..b)\gets A[m..r)$; mover $A[l..m)\to A[l{+}b..r)$ (de derecha a izquierda); $A[l..l{+}b)\gets\buf[0..b)$; \KwRet}
  \If(\tcp*[f]{puente}){$a > b$ \KwAnd $a-b \le \Bsz$}{
    $d \gets a-b$;\quad $\buf[0..d) \gets A[l+b..m)$\;
    \lFor{$i \gets 0$ \KwTo $b-1$}{$A[l+b+i] \gets A[l+i]$;\; $A[l+i] \gets A[l+b+d+i]$}
    $A[l+2b..r) \gets \buf[0..d)$; \KwRet\;
  }
  \If(\tcp*[f]{puente simétrico}){$b > a$ \KwAnd $b-a \le \Bsz$}{
    $d \gets b-a$;\quad $\buf[0..d) \gets A[m..m+d)$\;
    \lFor{$i \gets a-1$ \KwTo $0$}{$A[l+d+i] \gets A[m+d+i]$;\; $A[m+d+i] \gets A[l+i]$}
    $A[l..l+d) \gets \buf[0..d)$; \KwRet\;
  }
  \TRF{$l, m, r$}\;
}
\BlankLine
\Fn(\tcp*[h]{idea: \emph{trinity rotation} de blitsort}){\TRF{$l, m, r$}}{
  \lIf{$m-l = r-m$}{intercambiar $A[l..m)$ con $A[m..r)$; \KwRet}
  $p_a \gets l$;\; $p_b \gets m$;\; $p_c \gets m$;\; $p_d \gets r$\;
  \For{$k \gets 1$ \KwTo $\lfloor \min(m-l,\ r-m)/2 \rfloor$}{
    $t \gets A[{-}{-}p_b]$;\; $A[p_b] \gets A[p_a]$;\; $A[p_a{+}{+}] \gets A[p_c]$;\; $A[p_c{+}{+}] \gets A[{-}{-}p_d]$;\; $A[p_d] \gets t$\;
  }
  \eIf{$m-l < r-m$}{
    \For{$k \gets 1$ \KwTo $\lfloor (p_d-p_c)/2\rfloor$}{$t \gets A[p_c]$;\; $A[p_c{+}{+}] \gets A[{-}{-}p_d]$;\; $A[p_d] \gets A[p_a]$;\; $A[p_a{+}{+}] \gets t$\;}
  }{
    \For{$k \gets 1$ \KwTo $\lfloor (p_b-p_a)/2\rfloor$}{$t \gets A[{-}{-}p_b]$;\; $A[p_b] \gets A[p_a]$;\; $A[p_a{+}{+}] \gets A[{-}{-}p_d]$;\; $A[p_d] \gets t$\;}
  }
  invertir $A[p_a..p_d)$\;
}
\end{algorithm}

\begin{lema}[Rotar]\label{lem:rotar}
\[
\hoare{A[l..r)=X\cat Y,\ |X|=m-l}{\Rotar(l,m,r)}{A[l..r)=Y\cat X\ \wedge\ A\equiv A_1 \text{ fuera de } [l,r)}
\]
y escribe a lo sumo $2(r-l)$ elementos.
\end{lema}
\begin{proof}
Sean $a=|X|$, $b=|Y|$ y, sin pérdida de generalidad, $l=0$ y $N=a+b$. Cada lazo tiene
una cantidad fija de iteraciones, así que todas las ramas terminan.

\emph{Ramas con buffer.} Si $a\le\Bsz$, $X$ queda guardado en $\buf$. Mover $Y$ hacia
la izquierda, de izquierda a derecha, solo escribe en posiciones ya leídas, porque el
destino precede al origen. Al final $\buf$ se copia a $[b,N)$. La rama $b\le\Bsz$ es
simétrica: el movimiento hacia la derecha se hace de derecha a izquierda. Escrituras:
$2a+b\le 2N$.

\emph{Puente, $a>b$, $d=a-b\in[1,\Bsz]$.} Se escribe $X=A_1\cat A_2$ con $|A_1|=b$,
$|A_2|=d$, e $Y=B_b$ con $|B_b|=b$. El invariante del lazo, en la iteración $i$
($0\le i\le b$), es
\[
I(i):\ A[0..i)=B_b[0..i),\ A[i..b)=A_1[i..b),\ A[b..b{+}i)=A_1[0..i),\
A[b{+}d{+}i..N)=B_b[i..b),\ \buf[0..d)=A_2,
\]
y las posiciones $[b+i,\,b+d+i)$ (que son $d\ge 1$) están \emph{libres}: su contenido
no se usa. Para mostrar que la iteración $i$ preserva $I$: escribe en $A[b+i]$, que
está libre porque $b+i\in[b+i,b+d+i)$, el valor $A[i]=A_1[i]$. Luego escribe en
$A[i]$, que ya fue leída, el valor $A[b+d+i]=B_b[i]$. Las libres pasan a ser
$[b+i+1,\,b+d+i+1)$. La variante es $b-i$. En $i=b$:
$A[0..b)=B_b$, $A[b..2b)=A_1$, y la copia de $\buf$ a $[2b,N)$ deja
$B_b\cat A_1\cat A_2=Y\cat X$. Escrituras: $2d+2b=2a\le 2N$. El puente simétrico es
el mismo argumento reflejado, con el lazo descendente.

\emph{Triple reversión fusionada.} Si $a=b$, intercambiar los bloques es la rotación.
Si no, se definen
\[
M=\rev(X)\cat\rev(Y),\qquad F=\rev(M)=Y\cat X,\qquad\text{de modo que } F[q]=M[N-1-q].
\]
Explícitamente, $M[q]=X[a-1-q]$ para $q<a$ y $M[a+q]=Y[b-1-q]$ para $q<b$.

Caso $a<b$. Sea $K=\lfloor a/2\rfloor$. \textbf{Lazo 1}, invariante $J_1(k)$ para
$0\le k\le K$: las posiciones $[0,k)$ y $[N{-}k,N)$ valen $F$; $[a{-}k,a)$ y
$[a,a{+}k)$ valen $M$; $[k,a{-}k)$ y $[a{+}k,N{-}k)$ conservan $A_0$; y además
$p_a=k$, $p_b=a-k$, $p_c=a+k$, $p_d=N-k$. La iteración $k$ toca las posiciones
$k<a-1-k<a+k<N-1-k$, que son distintas porque $k<K$ y $a<b$. Hace
$A[a{-}1{-}k]\gets X[k]=M[a{-}1{-}k]$, $A[k]\gets Y[k]=F[k]$,
$A[a{+}k]\gets Y[b{-}1{-}k]=M[a{+}k]$ y $A[N{-}1{-}k]\gets X[a{-}1{-}k]=F[N{-}1{-}k]$.
Todos los valores leídos eran de $A_0$ por $J_1(k)$. Luego vale $J_1(k+1)$. Si $a$ es
impar, la posición central $K$ de $X$ cumple $A_0[K]=X[K]=M[K]$, así que al terminar
todo $[K,a+K)$ vale $M$.

\textbf{Lazo 2}, invariante $J_2(j)$: $[0,K{+}j)$ vale $F$; $[K{+}j,\,a{+}K{+}j)$
vale $M$; $[a{+}K{+}j,\,N{-}K{-}j)$ conserva $A_0$; y $[N{-}K{-}j,N)$ vale $F$. La
iteración toma $t=A[p_c]=Y[K{+}j]$, escribe $A[p_c]\gets A[p_d{-}1]=Y[b{-}1{-}K{-}j]
=M[a{+}K{+}j]$, $A[p_d{-}1]\gets A[p_a]=M[K{+}j]=F[N{-}1{-}K{-}j]$ y
$A[p_a]\gets t=Y[K{+}j]=F[K{+}j]$. Las tres posiciones son distintas mientras
$j<\lfloor(b-2K)/2\rfloor$. Al terminar, la zona que conserva $A_0$ tiene tamaño 0 o
1 y, si es 1, es el punto fijo de $\rev(Y)$, así que también vale $M$.

\textbf{Reversión final.} El centro $[p_a,p_d)=[K{+}J,\,N{-}K{-}J)$ es simétrico
respecto de $(N{-}1)/2$ y vale $M$. Invertirlo da
$A[q]=M[N{-}1{-}q]=F[q]$ en el centro. Por lo tanto $A=F=Y\cat X$.

El caso $a>b$ es simétrico: el lazo 2 opera sobre el resto de $X$, con
$J_2'(j)$: $[0,K{+}j)$ vale $F$, $[K{+}j,\,a{-}K{-}j)$ conserva $A_0$,
$[a{-}K{-}j,\,N{-}K{-}j)$ vale $M$ y $[N{-}K{-}j,N)$ vale $F$, con $K=\lfloor b/2\rfloor$.
Las lecturas de $A[p_d-1]$ caen dentro de la zona $M$ y nunca coinciden con $p_b$,
porque $b\ge 1$.

Escrituras: el lazo 1 hace 5 por cada 4 elementos, el lazo 2 hace 4 por cada 3 y la
reversión 3 por cada 2. En total, a lo sumo $1{,}5\,N+O(1)\le 2N$ para $N\ge 2$.
\end{proof}

%======================================================================
\section{Partición estable con memoria $O(1)$}

\begin{algorithm}[H]
\caption{Partición de bloque, partición estable y fusión}\label{alg:part}
\Fn(\tcp*[h]{$e-s\le\Bsz$}){\PartBloque{$s, e, P$}}{
  $m \gets 0$\;
  \For(\tcp*[f]{sin branch}){$i \gets 0$ \KwTo $e-s-1$}{
    $x \gets A[s+i]$;\; $t \gets P(x)$;\; $\buf[i-m] \gets x$;\; $A[s+m] \gets x$;\; $m \gets m + t$\;
  }
  $A[s+m..e) \gets \buf[0..e-s-m)$;\quad \KwRet $s+m$\;
}
\BlankLine
\Fn{\Part{$l, r, P$}}{
  pila $S \gets \emptyset$\tcp*{entradas $(\text{inicio}, \text{corte}, \text{nivel})$}
  \For{$s \gets l$ \KwTo $r-1$ \textbf{con paso} $\Bsz$}{
    $e\gets\min(s+\Bsz,\,r)$;\quad \textsf{push}$(S, (s,\ \PartBloque{$s, e, P$},\ 0))$\;
    \lWhile{$|S|\ge 2$ \KwAnd $\mathrm{nivel}(S_{-2}) = \mathrm{nivel}(S_{-1})$}{\Fusionar{$S$}}
  }
  \lWhile{$|S| \ge 2$}{\Fusionar{$S$}}
  \KwRet $\mathrm{corte}(S_0)$\;
}
\BlankLine
\Fn{\Fusionar{$S$}}{
  $(s_1, c_1, v_1) \gets S_{-2}$;\quad $(s_2, c_2, v_2) \gets S_{-1}$\;
  \Rotar{$c_1,\ s_2,\ c_2$};\quad $S_{-2} \gets (s_1,\ c_1 + c_2 - s_2,\ \max(v_1, v_2)+1)$;\quad \textsf{pop}$(S)$\;
}
\end{algorithm}

\begin{definicion}
Un segmento $[s,e)$ está \emph{$P$-partido con corte $c$} si
$A[s..e)=\mathrm{Part}_P(C)$, donde $C$ es su contenido al empezar la llamada a
\Part, y $c=s+|\filt(P,C)|$.
\end{definicion}

\begin{lema}[ParticiónBloque]\label{lem:partbloque}
\[
\hoare{e-s\le\Bsz,\ A[s..e)=C}{c\gets\PartBloque(s,e,P)}{A[s..e)=\mathrm{Part}_P(C)\ \wedge\ c=s+|\filt(P,C)|}
\]
\end{lema}
\begin{proof}
Invariante para $0\le i\le e-s$: $m=|\filt(P,C[0..i))|$,
$A[s..s{+}m)=\filt(P,C[0..i))$, $\buf[0..i{-}m)=\filt(\neg P,C[0..i))$ y
$A[s{+}i..e)=C[i..)$. La variante es $e-s-i$.

Las escrituras de la iteración $i$ caen en $A[s+m]$, con $s+m\le s+i$, una posición ya
leída, y en $\buf[i-m]$, con $i-m<\Bsz$. Si $t$ es verdadero, $A[s+m]=x$ queda dentro
del prefijo y $m$ crece; el valor escrito en $\buf[i-m]$ queda fuera de la zona útil,
porque $(i{+}1)-(m{+}1)=i-m$. Si $t$ es falso, $\buf[i-m]=x$ extiende los $\neg P$ y
la escritura en $A[s+m]$ queda fuera del prefijo útil.

Al terminar, la copia final coloca $\filt(\neg P,C)$ detrás del prefijo.
\end{proof}

\begin{lema}[Fusión]\label{lem:fusion}
Si $[s_1,s_2)$ está $P$-partido con corte $c_1$ y $[s_2,e)$ lo está con corte $c_2$,
después de $\Rotar(c_1,s_2,c_2)$ el segmento $[s_1,e)$ está $P$-partido con corte
$c_1+c_2-s_2$.
\end{lema}
\begin{proof}
Antes de rotar, $A[s_1..e)=T_1\cat F_1\cat T_2\cat F_2$. Rotar $F_1\cat T_2$ da
$T_1\cat T_2\cat F_1\cat F_2$, y $\filt$ de una concatenación es la concatenación de
los filtros. El nuevo corte es $s_1+|T_1|+|T_2|$.
\end{proof}

\begin{lema}[ParticiónEstable]\label{lem:partest}
Sea $L=r-l$ y $k=\lceil L/\Bsz\rceil$. Entonces
\[
\hoare{A[l..r)=C}{c\gets\Part(l,r,P)}{A[l..r)=\mathrm{Part}_P(C)\ \wedge\ c=l+|\filt(P,C)|}.
\]
Además, la pila nunca supera $\lfloor\log_2 k\rfloor+2$ entradas, y se hacen a lo sumo
$3L+4L\lfloor\log_2 k\rfloor$ escrituras y exactamente $L$ evaluaciones de $P$. Si $P$
es $\eqv$-invariante, la operación es OE.
\end{lema}
\begin{proof}
\emph{Invariante del lazo externo:} las entradas de la pila, de abajo hacia arriba,
describen segmentos contiguos $[s_0,s_1),\dots,[s_{t-1},s)$ con $s_0=l$, cada uno
$P$-partido con su corte. Además, la entrada de nivel $v$ cubre exactamente $2^v$
bloques, y los niveles son estrictamente decrecientes de abajo hacia arriba fuera del
\texttt{mientras} interno. Esto último es un contador binario: apilar suma 1 y cada
fusión de dos niveles iguales $v$ es un acarreo a $v+1$.

El Lema~\ref{lem:partbloque} da el bloque nuevo partido y el Lema~\ref{lem:fusion}
preserva el invariante en cada fusión. Variantes: $r-s$ para el lazo externo y $|S|$
para los lazos de fusión. Al final queda una sola entrada que cubre $[l,r)$.

La cantidad de entradas es la cantidad de unos del contador, a lo sumo
$\lfloor\log_2 k\rfloor+1$, más una transitoria. Como $n<2^{63}$ y $\Bsz=2^9$, son a
lo sumo 56 entradas; la implementación reserva 64.

\emph{Escrituras.} Cada elemento participa en a lo sumo $\lfloor\log_2 k\rfloor$
fusiones durante la fase de contador (cada una sube en 1 el nivel de su entrada) y en
a lo sumo $\lfloor\log_2 k\rfloor$ durante el colapso final. Una fusión rota un
subrango del segmento fusionado y, por el Lema~\ref{lem:rotar}, escribe a lo sumo
2 por elemento del segmento. Los bloques cuestan $\le 3L$ escrituras.

\emph{OE.} Es el Lema~\ref{lem:part}(a), porque el estado final es
$\mathrm{Part}_P(C)$.
\end{proof}

%======================================================================
\section{Merges}

\begin{algorithm}[H]
\caption{Merges con buffer y \emph{parity merge}}\label{alg:merges}
\Fn(\tcp*[h]{$m-l\le\Bsz$}){\MIzq{$l, m, r$}}{
  $a\gets m-l$;\quad $\buf[0..a)\gets A[l..m)$;\quad $p\gets 0$;\ $q\gets m$;\ $o\gets l$\;
  \While{$p<a$ \KwAnd $q<r$}{
    \leIf{$A[q]\lt\buf[p]$}{$A[o]\gets A[q]$; $q\gets q+1$}{$A[o]\gets\buf[p]$; $p\gets p+1$}
    $o\gets o+1$\;
  }
  $A[o..o+a-p)\gets\buf[p..a)$\;
}
\BlankLine
\Fn(\tcp*[h]{$r-m\le\Bsz$; simétrico, hacia atrás}){\MDer{$l, m, r$}}{
  igual que \MIzq con la derecha en $\buf$, recorriendo desde el final, y tomando la
  izquierda solo si $\buf[j-1]\lt A[i-1]$\;
}
\BlankLine
\Fn(\tcp*[h]{fuera de lugar: $D$ disjunto de $X$ e $Y$}){\Parity{$X[0..a),\ Y[0..b),\ D$}}{
  $h_x, h_y \gets 0$;\quad $t_x \gets a-1$;\quad $t_y \gets b-1$;\quad $h_d \gets 0$;\quad $t_d \gets a+b-1$\;
  \For{$k \gets 1$ \KwTo $\min(a,b)$}{
    $u \gets (Y[h_y] \lt X[h_x])$;\; $D[h_d{+}{+}] \gets u\,?\,Y[h_y]:X[h_x]$;\; $h_y \gets h_y{+}u$;\; $h_x \gets h_x{+}\neg u$\;
    $v \gets (Y[t_y] \lt X[t_x])$;\; $D[t_d{-}{-}] \gets v\,?\,X[t_x]:Y[t_y]$;\; $t_x \gets t_x{-}v$;\; $t_y \gets t_y{-}\neg v$\;
  }
  fusionar $X[h_x..t_x{+}1)$ e $Y[h_y..t_y{+}1)$ en $D[h_d..t_d{+}1)$ con el lazo clásico\;
}
\end{algorithm}

\begin{lema}[MergeIzq / MergeDer]\label{lem:mizq}
\[
\hoare{\sortd(X)\wedge\sortd(Y),\ A[l..r)=X\cat Y,\ |X|\le\Bsz}{\MIzq(l,m,r)}{A[l..r)=\Mrg(X,Y)}
\]
con a lo sumo $|X|+|Y|$ comparaciones. Lo mismo vale para \MDer con $|Y|\le\Bsz$.
\end{lema}
\begin{proof}
Invariante: $o=l+p+(q-m)$, $A[l..o)=\Mrg(X[0..p),Y[0..q{-}m))$,
$\buf[p..a)=X[p..)$ y $A[q..r)=Y[q{-}m..)$. Además $o\le q$, y la desigualdad es
estricta si $p<a$; esto asegura que no se sobrescribe ningún elemento de $Y$ sin leer.
Cada paso es el paso de cabeza (Lema~\ref{lem:merge}(a)). La variante es $(a-p)+(r-q)$.

Al salir: si $p=a$, el resto de $Y$ ya está en su lugar porque $o=q$; si $q=r$, la
copia final completa $X$. Ese resto es un segmento final de $\Mrg$ por el
Lema~\ref{lem:merge}(c). \MDer usa el paso de cola (Lema~\ref{lem:merge}(b)) con el
invariante simétrico.
\end{proof}

\begin{lema}[ParityMerge]\label{lem:parity}
\[
\hoare{\sortd(X[0..a))\wedge\sortd(Y[0..b)),\ D\cap(X\cup Y)=\emptyset}{\Parity(X,Y,D)}{D[0..a{+}b)=\Mrg(X,Y)}
\]
con a lo sumo $a+b$ comparaciones. Las primeras $\min(a,b)$ iteraciones no hacen
ningún chequeo de bordes.
\end{lema}
\begin{proof}
Sea $Z=\Mrg(X,Y)$, $N=a+b$ y $K=\min(a,b)$. La observación clave es que \emph{la
cabeza y la cola son dos merges clásicos independientes}, uno hacia adelante y otro
hacia atrás, que leen fuentes que no se modifican.

\emph{Cabeza.} Invariante tras $k$ iteraciones: $h_x+h_y=k$ y
$D[0..k)=Z[0..k)=\Mrg(X[0..h_x),Y[0..h_y))$. Para aplicar el
Lema~\ref{lem:merge}(a) hay que garantizar $h_x<a$ y $h_y<b$ antes de cada paso. Eso
vale porque $h_x,h_y\le k<K\le\min(a,b)$. \emph{Por eso no hacen falta chequeos de
bordes.}

\emph{Cola.} Simétricamente, $D[N{-}k..N)=Z[N{-}k..N)$ con $t_x=a-1-c_x\ge a-1-k\ge 0$
y lo mismo para $t_y$ (Lema~\ref{lem:merge}(b)).

\emph{Las dos partes no se pisan.} Como $2K\le N$, los segmentos $Z[0..K)$ y
$Z[N{-}K..N)$ son disjuntos. Como $\Mrg$ preserva el orden interno de $X$, la cabeza
consumió exactamente los primeros $h_x$ elementos de $X$ en $Z$ y la cola los últimos
$c_x$. Al ser disjuntos, $h_x\le a-c_x=t_x+1$, y lo mismo para $Y$.

\emph{Resto.} Por el Lema~\ref{lem:merge}(c),
$Z[K..N{-}K)=\Mrg(X[h_x..t_x{+}1),Y[h_y..t_y{+}1))$, que es exactamente lo que
calcula el lazo final, con la variante habitual.
\end{proof}

\begin{observacion}
La versión ``igual largo'' de quadsort es el caso $a=b$, donde el resto es vacío.
Usar $K=\min(a,b)$ en lugar de exigir largos casi iguales es la generalización que
emplean las hojas de tamaño arbitrario de SotSort.
\end{observacion}

%======================================================================
\section{Hojas}

\begin{algorithm}[H]
\caption{Inserción, ping-pong y hoja}\label{alg:hoja}
\Fn{\Insercion{$l, r$}}{
  \For{$i \gets l+1$ \KwTo $r-1$}{
    \If{$A[i] \lt A[i-1]$}{
      $v\gets A[i]$;\quad $j\gets i$\;
      \lWhile{$j>l$ \KwAnd $v\lt A[j-1]$}{$A[j]\gets A[j-1]$; $j\gets j-1$}
      $A[j]\gets v$\;
    }
  }
}
\BlankLine
\Fn(\tcp*[h]{$c=2^k$, $2\le c\le\Bsz$}){\PingPong{$P[0..c)$}}{
  $\Lambda\gets k-1$;\quad $\mathit{out}\gets(\Lambda \text{ impar})\;?\;\buf:P$\;
  \For{$i\gets 0,2,\dots,c-2$}{
    $x\gets P[i]$;\; $y\gets P[i{+}1]$;\; $s\gets(y\lt x)$;\; $\mathit{out}[i]\gets s\,?\,y:x$;\; $\mathit{out}[i{+}1]\gets s\,?\,x:y$\;
  }
  $\mathit{src}\gets\mathit{out}$;\quad $\mathit{dst}\gets(\mathit{out}=P)\,?\,\buf:P$\;
  \For{$w\gets 2,4,\dots,c/2$}{
    \lFor{$i\gets 0,\,2w,\,4w,\dots$}{\Parity{$\mathit{src}[i..i{+}w),\ \mathit{src}[i{+}w..i{+}2w),\ \mathit{dst}{+}i$}}
    intercambiar $\mathit{src}$ y $\mathit{dst}$\;
  }
}
\BlankLine
\Fn(\tcp*[h]{$r-l\le\Bsz$}){\Hoja{$l, r$}}{
  $e\gets r$;\quad $q\gets r-l$;\quad $\mathit{primero}\gets\mathbf{true}$\;
  \While{$q>0$}{
    $c\gets q \mathbin{\&} (-q)$;\quad $s\gets e-c$\tcp*{trozo más chico que queda}
    \lIf{$c>1$}{\PingPong{$A[s..s{+}c)$}}
    \If{\KwNot $\mathit{primero}$ \KwAnd $A[s+c]\lt A[s+c-1]$}{
      \Parity{$A[s..s{+}c),\ A[s{+}c..r),\ \buf$};\quad $A[s..r)\gets\buf[0..r{-}s)$\;
    }
    $\mathit{primero}\gets\mathbf{false}$;\quad $e\gets s$;\quad $q\gets q-c$\;
  }
}
\end{algorithm}

\begin{lema}[Inserción]\label{lem:insercion}
$\hoare{A[l..r)=C}{\Insercion(l,r)}{\sortd(A[l..r))\ \wedge\ \text{la transformación es OE}}$.
\end{lema}
\begin{proof}
El invariante del lazo externo es que $A[l..i)$ es una disposición ordenada OE de
$C[0..i{-}l)$. El elemento $v$ solo se desplaza por encima de elementos
$A[j-1]$ con $v\lt A[j-1]$, que no son equivalentes a $v$. Se detiene ante el primero
con $v\nlt A[j-1]$, y ese elemento queda antes de $v$, de modo que el orden entre
equivalentes se preserva. Variantes: $r-i$ para el lazo externo y $j-l$ para el
interno.
\end{proof}

\begin{lema}[PingPongPot2]\label{lem:pingpong}
Con $c=2^k\le\Bsz$, $\PingPong(P)$ deja $P$ ordenado mediante una transformación OE,
con a lo sumo $c\,k$ comparaciones.
\end{lema}
\begin{proof}
Paso de pares: se intercambian solo si $y\lt x$, así que dos equivalentes conservan
su orden. Invariante del lazo de niveles, al entrar con ancho $w$: cada bloque
alineado $\mathit{src}[jw..(j{+}1)w)$ está ordenado y contiene exactamente los
elementos de $P_0[jw..(j{+}1)w)$, y $\mathit{src}$ es una disposición OE de $P_0$.

Cada \Parity cumple su precondición, porque $\mathit{dst}$ es disjunto de
$\mathit{src}$, y produce $\Mrg$ de dos bloques adyacentes. Eso es OE por el
Lema~\ref{lem:merge}(d), y el invariante vale para $2w$.

Paridad: hay $\Lambda=k-1$ niveles y cada uno alterna el destino. Si el paso de pares
escribió en $\buf$ exactamente cuando $\Lambda$ es impar, al terminar se cumple
$\mathit{src}=P$.
\end{proof}

\begin{lema}[Hoja]\label{lem:hoja}
$\hoare{r-l\le\Bsz,\ A[l..r)=C}{\Hoja(l,r)}{\sortd(A[l..r))\ \wedge\ \text{OE}}$, con
$O((r-l)\log(r-l))$ comparaciones.
\end{lema}
\begin{proof}
Invariante: $A[e..r)$ es una disposición ordenada OE de $C[e{-}l..)$, y
$q=e-l$. El trozo $[s,e)$ tiene tamaño potencia de 2 (el bit más bajo de $q$) y queda
ordenado por el Lema~\ref{lem:pingpong}.

Si $A[s+c]\nlt A[s+c-1]$, la concatenación ya está ordenada y es trivialmente OE. Si
no, \Parity escribe en $\buf$, que es disjunto y tiene capacidad suficiente porque
$r-s\le\Bsz$, el valor $\Mrg$ del trozo (a la izquierda) y del sufijo. Eso es OE por
el Lema~\ref{lem:merge}(d), y la copia lo devuelve a su lugar. La variante es $q$,
que pierde un bit en cada iteración.
\end{proof}

%======================================================================
\section{SymMerge iterativo y merge sort de respaldo}

\begin{algorithm}[H]
\caption{SymMerge iterativo y merge sort de respaldo}\label{alg:sym}
\Fn{\SymMerge{$l_0, m_0, r_0$}}{
  pila $W\gets[(l_0,m_0,r_0)]$\;
  \While{$W\ne\emptyset$}{
    $(l,m,r)\gets$ \textsf{pop}$(W)$\;
    \lIf{$l\ge m$ \KwOr $m\ge r$ \KwOr $A[m]\nlt A[m-1]$}{\KwCont}
    $l\gets$ primer $i\in[l,m)$ con $A[m]\lt A[i]$;\quad $r\gets$ primer $j\in[m,r)$ con $A[j]\nlt A[m-1]$\;
    $n_l\gets m-l$;\quad $n_r\gets r-m$\;
    \lIf{$n_l+n_r\le\Bsz$}{$\buf\gets A[l..r)$; \Parity{$\buf[0..n_l),\ \buf[n_l..n_l{+}n_r),\ A{+}l$}; \KwCont}
    \lIf{$\min(n_l,n_r)\le\Bsz$}{\MIzq{$l,m,r$} o \MDer{$l,m,r$} según el lado chico; \KwCont}
    \eIf{$n_l\ge n_r$}{
      $m_1\gets l+\lfloor n_l/2\rfloor$;\quad $m_2\gets$ primer $j\in[m,r)$ con $A[j]\nlt A[m_1]$\;
    }{
      $m_2\gets m+\lfloor n_r/2\rfloor$;\quad $m_1\gets$ primer $i\in[l,m)$ con $A[m_2]\lt A[i]$\;
    }
    \Rotar{$m_1, m, m_2$};\quad $\nu\gets m_1+(m_2-m)$\;
    apilar $(l,m_1,\nu)$ y $(\nu,m_2,r)$, \emph{el más grande primero}\;
  }
}
\BlankLine
\Fn{\Respaldo{$l, r$}}{
  \lFor{$s\gets l,\,l{+}24,\dots$}{\Insercion{$s,\ \min(s{+}24,r)$}}
  \For{$w\gets 24, 48, 96,\dots$ \KwAnd $w<r-l$}{
    \lFor{$x\gets l,\,l{+}2w,\dots$ \KwAnd $x+w<r$}{\SymMerge{$x,\ x{+}w,\ \min(x{+}2w,\,r)$}}
  }
}
\end{algorithm}

\begin{lema}[Recorte]\label{lem:recorte}
Sea $X\cat Y=A[l..r)$ con $X=A[l..m)$, $Y=A[m..r)$ ordenadas, y sea $Z=\Mrg(X,Y)$.
\begin{enumerate}[label=(\alph*),nosep]
\item Si $A[m]\nlt A[m-1]$, entonces $X\cat Y=Z$.
\item Con $l'$ y $r'$ como en la línea 5, $Z = A[l..l')\cat\Mrg(A[l'..m),A[m..r'))\cat A[r'..r)$.
\end{enumerate}
\end{lema}
\begin{proof}
(a) Todo $x\in X$ cumple $x\les X_{\max}=A[m-1]\les A[m]\les y$ para todo $y\in Y$
(Lema~\ref{lem:swo}(d)). Ante equivalencias, $X$ va primero; por lo tanto la
concatenación cumple la Definición~\ref{def:merge}.
(b) Los elementos de $A[l..l')$ cumplen $x\les A[m]$, el mínimo de $Y$. Así preceden
a todo $Y$, incluso con equivalencia, porque son de $X$. Además preceden al resto de
$X$ por el orden de $X$. Simétricamente, los de $A[r'..r)$ cumplen $A[m-1]\les y$,
así que van después de todo $X$, y después del resto de $Y$ por el orden de $Y$.
Luego son prefijo y sufijo de $Z$, y el centro es el $\Mrg$ de lo que queda
(Lema~\ref{lem:merge}(c)).
\end{proof}

\begin{lema}[División de SymMerge]\label{lem:symsplit}
En la rama $n_l\ge n_r$, después de \Rotar se cumple
\[
\Mrg(X,Y)=\Mrg(X[l..m_1),Y[m..m_2))\cat\Mrg(X[m_1..m),Y[m_2..r)),
\]
y la rotación es OE. Lo mismo vale en la otra rama.
\end{lema}
\begin{proof}
Rama $n_l\ge n_r$. Se cumple $Y[m..m_2)\lt A[m_1]$ y $A[m_1]\les Y[m_2..r)$.
\begin{itemize}[nosep]
\item Todo $y\in Y[m..m_2)$ cumple $y\lt A[m_1]\les x'$ para todo
  $x'\in X[m_1..m)$, así que $y$ precede estrictamente a $x'$.
\item Todo $x\in X[l..m_1)$ cumple $x\les A[m_1]\les y'$ para todo
  $y'\in Y[m_2..r)$, así que $x$ precede a $y'$, y ante equivalencia también, porque
  $x$ es de $X$.
\end{itemize}
Entonces, en $Z$, todo el grupo izquierdo precede al grupo derecho, y por el
Lema~\ref{lem:merge}(c) cada grupo es el $\Mrg$ de sus partes. La rotación intercambia
$X[m_1..m)$ con $Y[m..m_2)$. Entre esos dos bloques no hay pares equivalentes, porque
$y\lt A[m_1]\les x'$; por el Lema~\ref{lem:rotsegura} es OE.

En la otra rama, $X[l..m_1)\les A[m_2]\lt X[m_1..m)$ y el argumento es simétrico.
\end{proof}

\begin{lema}[SymMerge]\label{lem:sym}
Sea $s=r_0-l_0$. Entonces
\[
\hoare{\sortd(A[l_0..m_0))\wedge\sortd(A[m_0..r_0))}{\SymMerge(l_0,m_0,r_0)}{A[l_0..r_0)=\Mrg(\text{ambas})}
\]
y:
\begin{enumerate}[label=(\alph*),nosep]
\item la pila $W$ nunca supera $\lfloor\log_2 s\rfloor+2$ entradas;
\item la profundidad del árbol de subproblemas es a lo sumo
  $\log_{9/7}(s/2\Bsz)+2$;
\item hace $O(s\log(s/\Bsz))$ escrituras;
\item hace a lo sumo $s+\frac{s}{2\Bsz}\,D\,(3\lceil\log_2(s{+}1)\rceil+1)$
  comparaciones, donde $D$ es la profundidad de (b).
\end{enumerate}
\end{lema}
\begin{proof}
\emph{Invariante:} los subproblemas $(l,m,r)$ que están en $W$ cubren intervalos
disjuntos. En cada uno, $A[l..m)$ y $A[m..r)$ están ordenados. Además, si cada
subproblema pendiente se reemplazara por su $\Mrg$, el arreglo completo en
$[l_0,r_0)$ valdría $\Mrg$ de la entrada.

Las ramas que terminan un subproblema son correctas por los
Lemas~\ref{lem:recorte}, \ref{lem:parity} y~\ref{lem:mizq}. La rama de división es
correcta por el Lema~\ref{lem:symsplit}.

\emph{Terminación.} En una división, $n_l,n_r>\Bsz\ge 1$, así que
$\lfloor n_l/2\rfloor\ge 1$. Los dos hijos son no vacíos y cada uno es estrictamente
más chico que el padre. El multiconjunto de tamaños pendientes decrece en el orden de
multiconjuntos de Dershowitz--Manna, que es bien fundado.

(a) Invariante adicional: si $I_1,\dots,I_t$ son las entradas de abajo hacia arriba,
$\sum_{j>i}|I_j|\le|I_i|$. Se mantiene porque los hijos suman a lo sumo el tamaño del
padre y el chico queda arriba. Entonces la suma desde $I_i$ hasta el tope al menos se
reduce a la mitad en cada paso, y $t\le\lfloor\log_2 s\rfloor+1$.

(b) Un nodo interno de tamaño $\sigma$ cumple $\sigma>2\Bsz$. Si $n_l\ge n_r$,
entonces $n_l\ge\sigma/2$ y cada hijo tiene a lo sumo
$\sigma-\lfloor n_l/2\rfloor\le\frac34\sigma+\frac12\le\frac79\sigma$ elementos,
porque $\sigma\ge 18$.

(c) Los nodos de una misma profundidad son disjuntos. Cada nodo interno escribe a lo
sumo $2\sigma$ (Lema~\ref{lem:rotar}) y cada hoja a lo sumo $3\sigma$.

(d) Las hojas comparan a lo sumo una vez por elemento y son disjuntas, lo que da $s$.
Cada nodo interno hace tres búsquedas binarias más una comparación inicial, y hay a lo
sumo $s/2\Bsz$ nodos internos por nivel.
\end{proof}

\begin{lema}[MergeRespaldo]\label{lem:respaldo}
\[\hoare{A[l..r)=C}{\Respaldo(l,r)}{\sortd(A[l..r))\ \wedge\ \text{OE}},\]
con $O(L\log L\log(L/\Bsz))$ escrituras y
$O(L\log L+\frac{L}{\Bsz}\log^3 L)$ comparaciones, donde $L=r-l$.
\end{lema}
\begin{proof}
Invariante del lazo en $w$: los bloques alineados de ancho $w$ están ordenados, y el
estado es una transformación OE de $C$ que no mueve elementos entre bloques. La
inserción lo establece para $w=24$ (Lema~\ref{lem:insercion}). Cada SymMerge sobre dos
bloques adyacentes lo extiende a $2w$ (Lemas~\ref{lem:sym} y~\ref{lem:merge}(d)).
Hay $\lceil\log_2(L/24)\rceil$ niveles, y las cotas salen de sumar el
Lema~\ref{lem:sym} sobre cada nivel.
\end{proof}

%======================================================================
\section{QuickEstable}

\begin{algorithm}[H]
\caption{\textsc{QuickEstable} (v6.1) y \textsc{Pivote}}\label{alg:quick}
\Fn{\Quick{$l_0, r_0$}}{
  pila $S \gets [\,(l_0, r_0, \beta_0, \bot)\,]$ con $\beta_0=\lfloor\log_2(r_0-l_0)\rfloor+1$\;
  \While{$S \neq \emptyset$}{
    $(l, r, \beta, \mathit{lb}) \gets$ \textsf{pop}$(S)$\;
    \While{$\mathbf{true}$}{
      $\lambda\gets r-l$\;
      \lIf{$\lambda \le \leaf$}{\Hoja{$l, r$}; \KwBreak}
      \lIf{$\beta \le 0$}{\lIf{\KwNot \Ordenado{$l,r$}}{\Respaldo{$l,r$}} \KwBreak}
      $p \gets$ \Pivote{$l, \lambda$}\tcp*{copia del valor de un elemento de $[l,r)$}
      \If{$\mathit{lb} \ne \bot$ \KwAnd $\mathit{lb} \nlt p$}{
        $l' \gets$ \Part{$l, r,\ x \mapsto p\nlt x$}\;
        \lIf{$l'-l < \lambda/16$}{$\beta\gets\beta-1$}\tcp*{\textbf{nuevo en v6.1}}
        $l\gets l'$; \KwCont\;
      }
      $m \gets$ \Part{$l, r,\ x \mapsto x \lt p$}\;
      \If{$m = l$}{
        $l' \gets$ \Part{$l, r,\ x \mapsto p\nlt x$}\;
        \lIf{$l'-l < \lambda/16$}{$\beta\gets\beta-1$}\tcp*{\textbf{nuevo en v6.1}}
        $l\gets l'$; \KwCont\;
      }
      \lIf{$\min(m-l,\ r-m) < \lambda/16$}{$\beta \gets \beta-1$}
      \leIf{$m - l < r - m$}{\textsf{push}$(S, (m, r, \beta, p))$; $r \gets m$}{\textsf{push}$(S, (l, m, \beta, \mathit{lb}))$; $l \gets m$; $\mathit{lb} \gets p$}
    }
  }
}
\BlankLine
\Fn{\Pivote{$l, \lambda$}}{
  \leIf{$\lambda \le N_9$}{\KwRet ninther: mediana de 3 medianas de 3 en 9 posiciones distintas de $[l,l{+}\lambda)$}{
    copiar $s=\min(\Bsz{-}1,\,2\lfloor\lambda^{1/3}\rfloor)$ (impar) muestras de $[l,l{+}\lambda)$ a $\buf$; \KwRet la mediana de $\buf[0..s)$}
}
\end{algorithm}

\begin{lema}[Pivote]\label{lem:pivote}
\Pivote no modifica $A$ y devuelve una copia del valor de algún elemento de
$A[l..l{+}\lambda)$.
\end{lema}
\begin{proof}
La mediana de 3 devuelve uno de sus 3 argumentos. Las muestras son copias de
posiciones $l+i\cdot\mathit{step}+\mathit{step}/2<l+\lambda$, y la mediana de esas
copias es una de ellas.
\end{proof}

\begin{definicion}[Estado de \Quick]\label{def:invq}
Sea $A_1$ el arreglo al llamar a $\Quick(l_0,r_0)$, y supongamos que $A_1[l_0..r_0)$
es coherente y $\elems(A_1[l_0..r_0))=\elems(\Omega[l_0..r_0))$. Un intervalo $I$ es
\begin{itemize}[nosep]
\item \emph{terminado} si $A[I]=\Omega[I]$;
\item \emph{pendiente} si $\elems(A[I])=\elems(\Omega[I])$ y $\coh(A[I])$.
\end{itemize}
El invariante $\mathit{INV}$ dice que $[l_0,r_0)$ se descompone en intervalos
disjuntos, cada uno terminado o pendiente; que los pendientes son exactamente los de
la pila más, dentro del lazo interno, el actual $[l,r)$; y que, si una entrada tiene
$\mathit{lb}\ne\bot$, todo $x$ de su intervalo cumple $x\nlt\mathit{lb}$.
\end{definicion}

\begin{teorema}[QuickEstable]\label{thm:quick}
\[
\hoare{\coh(A[l_0..r_0))\ \wedge\ \elems(A[l_0..r_0))=\elems(\Omega[l_0..r_0))}{\Quick(l_0,r_0)}{A[l_0..r_0)=\Omega[l_0..r_0)}.
\]
\end{teorema}
\begin{proof}
$\mathit{INV}$ vale al inicio: un único pendiente $[l_0,r_0)$. Se verifica que cada
rama del lazo interno lo preserva. Sea $I=[l,r)$ el intervalo actual, pendiente.

\emph{Hoja} (línea 7). Por el Lema~\ref{lem:hoja} queda $\sortd(A[I])$ y, como la
operación es OE (Lema~\ref{lem:oe}(a)), $A[I]$ sigue siendo coherente. Por el
Lema~\ref{lem:coh}, $A[I]=\Omega[I]$: el intervalo pasa a terminado.

\emph{Respaldo} (línea 8). Si \Ordenado devuelve verdadero, $A[I]$ ya está ordenado,
o era estrictamente descendente y se invirtió (Lema~\ref{lem:yaord}). En otro caso,
el Lema~\ref{lem:respaldo} da un resultado ordenado y OE. En ambos casos el
Lema~\ref{lem:coh} lo vuelve terminado.

\emph{Partición por $P_<$} (línea 15, con $m>l$). Por el Lema~\ref{lem:partest},
$A[I]=\filt(P_<,C)\cat\filt(\neg P_<,C)$, y la operación es OE porque $P_<$ es
$\eqv$-invariante (Lema~\ref{lem:part}(b)). Por el Lema~\ref{lem:part}(c) aplicado a
$\Omega[I]$, que está ordenado, $\Omega[I]$ también se parte en $l..m$ y $m..r$ con
los mismos conjuntos. Luego $[l,m)$ y $[m,r)$ son pendientes. Todo $x$ de $[m,r)$
cumple $x\nlt p$, lo que justifica $\mathit{lb}=p$ para ese hijo; el hijo izquierdo
hereda su $\mathit{lb}$ porque es subconjunto.

\emph{Separación de equivalentes} (líneas 10--13 y 17--19). Hay que ver que ningún
$x\in A[I]$ cumple $x\lt p$.
\begin{itemize}[nosep]
\item Rama de la línea 10. Si $x\lt p$, la transitividad negativa da
  $x\lt\mathit{lb}\vee\mathit{lb}\lt p$. Lo segundo contradice la guarda y lo primero
  contradice el invariante sobre $\mathit{lb}$.
\item Rama de la línea 17. Es directo: $m=l$ significa que ningún elemento cumple
  $P_<$.
\end{itemize}
Entonces $P_\le(x)\iff x\eqv p$, y el bloque $[l,l')$ queda formado exactamente por
los equivalentes a $p$, en orden coherente. Un bloque de equivalentes coherente está
ordenado y es coherente, y por el mismo argumento de partición de $\Omega[I]$ es
$\Omega[l..l')$: queda terminado. $[l',r)$ sigue pendiente.

Al vaciarse la pila no quedan pendientes, así que $A[l_0..r_0)=\Omega[l_0..r_0)$. La
terminación se prueba en el Lema~\ref{lem:prof}.
\end{proof}

\begin{lema}[Progreso, profundidad y pila]\label{lem:prof}
Sea $\lambda_0=r_0-l_0$.
\begin{enumerate}[label=(\alph*),nosep]
\item Cada iteración del lazo interno termina el intervalo actual o lo achica
  estrictamente; en una partición por $P_<$, los dos hijos son no vacíos.
\item Toda cadena de intervalos anidados procesados tiene longitud a lo sumo
  $D_Q=\beta_0+\lceil\log_{16/15}(\lambda_0/\leaf)\rceil$, es decir, $O(\log\lambda_0)$.
\item La pila $S$ nunca supera $\lfloor\log_2\lambda_0\rfloor+1\le 64$ entradas.
\item \Quick termina.
\end{enumerate}
\end{lema}
\begin{proof}
(a) Por el Lema~\ref{lem:pivote}, el elemento cuyo valor es $p$ está en $I$ y no
cumple $p\lt p$, así que va a la parte derecha: $m<r$. La rama de partición exige
$m>l$. En la separación de equivalentes, ese mismo elemento cumple $P_\le$, así que
$l'>l$.

(b) Cada paso de la cadena achica el tamaño a lo sumo por un factor $15/16$ (paso
``bueno'') o decrementa $\beta$ (paso ``malo''); esto vale tanto para las particiones
como para las separaciones, gracias al cambio de v6.1. Como $\beta$ se hereda y la
cadena corta cuando $\beta\le 0$, hay a lo sumo $\beta_0$ pasos malos. Los buenos
reducen el tamaño hasta $\leaf$ en a lo sumo $\lceil\log_{16/15}(\lambda_0/\leaf)\rceil$
pasos.

\emph{Sin el cambio de v6.1}, una separación que quita un solo elemento no consume
$\beta$, y la cota no se puede probar.

(c) Se apila el hijo grande y se continúa con el chico. Con $C$ el intervalo actual,
el invariante $\sum_{j>i}|I_j|+|C|\le|I_i|$ se mantiene igual que en el
Lema~\ref{lem:sym}(a).

(d) Se sigue de (a) y (b): hay finitos intervalos anidados y cada uno se procesa una
cantidad finita de veces.
\end{proof}

%======================================================================
\section{Camino adaptativo: detección y runs naturales}

\begin{algorithm}[H]
\caption{Detección y punto de entrada}\label{alg:main}
\Fn{\Ordenado{$l, r$}}{
  $i\gets$ \FinAsc{$l, r$}\tcp*{bloques de 32 sin branch; ver texto}
  \lIf{$i = r$}{\KwRet $\mathbf{true}$}
  $e_0\gets i$\tcp*{fin del prefijo no decreciente, se reutiliza}
  \If{$i = l+1$ \KwAnd \FinDesc{$l, r$} $= r$}{invertir $A[l..r)$; \KwRet $\mathbf{true}$}
  \KwRet $\mathbf{false}$\;
}
\BlankLine
\Fn(\tcp*[h]{heurística; no modifica $A$}){\PocosGiros{$n$}}{
  \lIf{$n\le 2^{16}$}{\KwRet $8\cdot\mathrm{giros}(2,n) < n$}
  \KwRet $8\cdot\sum_{c=0}^{63}\mathrm{giros}(b_c,\,b_c{+}256) < 64\cdot256$\tcp*{64 tramos equiespaciados}
}
\BlankLine
\Fn{\SotSort{$A, n$}}{
  \lIf{$n < 2$ \KwOr \Ordenado{$0, n$}}{\KwRet}
  \lIf{\PocosGiros{$n$}}{$\mathit{adaptativo}\gets\mathbf{true}$; \Natural{$0, n, e_0$}; \KwRet}
  \Quick{$0, n$}\;
}
\end{algorithm}

Aquí $\FinAsc(l,r)$ es el primer $j>l$ con $A[j]\lt A[j-1]$ (o $r$ si no hay), y
$\FinDesc(l,r)$ es el primer $j>l$ con $A[j]\nlt A[j-1]$. Ambos se calculan en
bloques de 32 pares, acumulando las comparaciones con un \textsc{or} en un entero de
64 bits, y después elemento por elemento. Por eso el resultado es exactamente la
definición. Se define $\mathrm{giros}(b,e)=\#\{i\in[b,e):\ [A[i]\lt A[i-1]]\ne
[A[i-1]\lt A[i-2]]\}$.

\begin{lema}[YaOrdenado]\label{lem:yaord}
Si \Ordenado devuelve verdadero, $A[l..r)$ queda ordenado y la transformación es OE.
Si devuelve falso, $A$ no cambia. En ambos casos $A[l..e_0)$ es no decreciente.
\end{lema}
\begin{proof}
Si $\FinAsc=r$, ningún par adyacente desciende y, por la transitividad de $\les$, todo
el segmento está ordenado. En la rama descendente, todo par adyacente cumple
$A[j]\lt A[j-1]$; por transitividad, todo par $i<j$ cumple $A[j]\lt A[i]$, así que no
hay equivalentes. La inversión da un segmento ordenado y es OE (vacuamente).
\end{proof}

\begin{observacion}[La heurística no afecta la corrección]\label{obs:heur}
\PocosGiros solo elige entre \Natural y \Quick. Ambos caminos son correctos para
cualquier entrada (Teoremas~\ref{thm:quick} y~\ref{thm:natural}). Por lo tanto el
muestreo puede equivocarse \emph{solo en rendimiento}.
\end{observacion}

\begin{algorithm}[H]
\caption{Merge por runs naturales con criterio de powersort}\label{alg:natural}
\Fn{\Potencia{$s_1, n_1, n_2, n$}}{
  $a\gets 2s_1+n_1$;\quad $b\gets a+n_1+n_2$;\quad $\rho\gets 0$\;
  \While{$\mathbf{true}$}{
    $\rho\gets\rho+1$\;
    \lIf{$a\ge n$}{$a\gets a-n$; $b\gets b-n$}
    \lElseIf{$b\ge n$}{\KwRet $\rho$}
    $a\gets 2a$;\quad $b\gets 2b$\;
  }
}
\BlankLine
\Fn{\Natural{$lo, hi, e_0$}}{
  pila $R\gets\emptyset$\tcp*{entradas $(s,\ \mathit{len},\ \mathit{pot})$}
  $i\gets lo$\;
  \While{$i<hi$}{
    \uIf{$i=lo$ \KwAnd $e_0>lo+1$}{$j\gets e_0$}
    \uElseIf{$i+1<hi$ \KwAnd $A[i+1]\lt A[i]$}{$j\gets$ \FinDesc{$i, hi$}; invertir $A[i..j)$}
    \lElse{$j\gets$ \FinAsc{$i, hi$}}
    \lIf{$j-i<32$}{$j\gets\min(hi,\,i{+}32)$; \Insercion{$i, j$}}
    \If{$R\ne\emptyset$}{
      $\rho\gets$ \Potencia{$s(R_{-1})-lo,\ \mathit{len}(R_{-1}),\ j-i,\ hi-lo$}\;
      \lWhile{$|R|\ge 2$ \KwAnd $\mathit{pot}(R_{-2})>\rho$}{\FusTope{$R$}}
      $\mathit{pot}(R_{-1})\gets\rho$\;
    }
    \textsf{push}$(R,(i,\ j-i,\ 0))$;\quad $i\gets j$\;
  }
  \lWhile{$|R|\ge 2$}{\FusTope{$R$}}
}
\BlankLine
\Fn{\FusTope{$R$}}{
  $(s_1,n_1,\cdot)\gets R_{-2}$;\quad $(s_2,n_2,\cdot)\gets R_{-1}$\;
  \SymMerge{$s_1,\ s_1{+}n_1,\ s_2{+}n_2$};\quad $\mathit{len}(R_{-2})\gets n_1+n_2$;\quad \textsf{pop}$(R)$\;
}
\end{algorithm}

\begin{teorema}[MergeNatural]\label{thm:natural}
Si $A[lo..e_0)$ es no decreciente, entonces
\[\hoare{A[lo..hi)=C}{\Natural(lo,hi,e_0)}{\sortd(A[lo..hi))\ \wedge\ \text{OE}}.\]
\end{teorema}
\begin{proof}
\emph{Invariante del lazo principal:} las entradas de $R$ describen segmentos
contiguos que cubren $[lo,i)$. Cada uno está ordenado y contiene exactamente los
elementos de $C$ en sus posiciones, y la transformación hecha hasta ahora es OE.
$A[i..hi)$ no fue tocado.

Detección de run: un run no decreciente ya está ordenado; uno estrictamente
descendente no tiene equivalentes y su inversión es OE (Lema~\ref{lem:yaord}); la
extensión por inserción es OE (Lema~\ref{lem:insercion}). \FusTope combina dos runs
adyacentes con SymMerge, lo que es correcto y OE por los Lemas~\ref{lem:sym}
y~\ref{lem:merge}(d).

Variantes: $hi-i$ crece en cada vuelta, porque $j>i$; $|R|$ decrece en cada fusión. Al
final queda un solo run que cubre $[lo,hi)$.
\end{proof}

\begin{lema}[Potencias]\label{lem:pot}
Sean dos runs adyacentes $[s_1,s_1{+}n_1)$ y $[s_2,s_2{+}n_2)$ con $s_2=s_1+n_1$, y
sean sus puntos medios normalizados $\alpha=\frac{2s_1+n_1}{2n}$ y
$\beta=\frac{2s_2+n_2}{2n}$. Entonces:
\begin{enumerate}[label=(\alph*),nosep]
\item $\Potencia$ termina y devuelve
  $\pi=\min\{\ell\ge 1:\ (\alpha,\beta] \text{ contiene un múltiplo de } 2^{-\ell}\}$;
\item $1\le\pi\le P_{\max}:=\lceil\log_2 n\rceil$, y no hay desbordamiento si
  $n<2^{63}$;
\item \emph{(potencias iguales)} si dos fronteras $x<y$ tienen potencia $\ell$, alguna
  frontera estrictamente entre ellas tiene potencia menor que $\ell$.
\end{enumerate}
\end{lema}
\begin{proof}
(a) Se tiene $0\le\alpha<\beta<1$, porque $2s_2+n_2<2n$. El lazo calcula los dígitos
binarios de $\alpha=a/2n$ y $\beta=b/2n$ de a uno: la condición $a\ge n$ equivale a
que el dígito actual de $\alpha$ sea 1. Devuelve la posición del primer dígito en que
difieren, y difieren en ese dígito si y solo si
$\lfloor 2^\ell\alpha\rfloor<\lfloor2^\ell\beta\rfloor$, es decir, si hay un múltiplo
de $2^{-\ell}$ en $(\alpha,\beta]$.

(b) Si coinciden en los primeros $\pi-1$ dígitos, $\beta-\alpha<2^{1-\pi}$. Como
$\beta-\alpha=\frac{n_1+n_2}{2n}\ge\frac1n$, se obtiene $2^{\pi-1}<n$. Además,
siempre $a,b<2n$ antes de duplicar, así que no hay desbordamiento.

(c) Sean $d_x\in(\alpha_x,\beta_x]$ y $d_y\in(\alpha_y,\beta_y]$ múltiplos de
$2^{-\ell}$. Ningún intervalo contiene un múltiplo de $2^{-(\ell-1)}$, así que $d_x$
y $d_y$ son múltiplos \emph{impares} de $2^{-\ell}$. Como $\beta_x\le\alpha_y$, se
tiene $d_x<d_y$, y $d=d_x+2^{-\ell}$ es múltiplo de $2^{-(\ell-1)}$ con
$d_x<d<d_y$.

Los intervalos $(\mu_t,\mu_{t+1}]$ entre puntos medios consecutivos, para las
fronteras $t$ de $x$ a $y$, cubren $(\alpha_x,\beta_y]$. El intervalo que contiene a
$d$ corresponde a una frontera de potencia $\le\ell-1$, y esa frontera no es ni $x$
ni $y$.
\end{proof}

\begin{lema}[Pila de runs y merges por elemento]\label{lem:natcost}
(a) Fuera del lazo de fusiones, las potencias de $R$ son estrictamente crecientes de
abajo hacia arriba; por lo tanto $|R|\le P_{\max}+1\le 64$.
(b) Cada elemento participa en a lo sumo $P_{\max}$ fusiones.
(c) \Natural hace $O(n\log n\log(n/\Bsz))$ escrituras y
$O(n\log n+\frac{n}{\Bsz}\log^3 n)$ comparaciones.
\end{lema}
\begin{proof}
(a) Tras el lazo de fusiones, $\mathit{pot}(R_{-2})\le\rho$. Supongamos que hay
igualdad. La frontera guardada en $R_{-2}$, llamémosla $x$, y la nueva, $z$, tienen
potencia $\rho$. Por el Lema~\ref{lem:pot}(c), hay una frontera $y$ estrictamente
entre ellas con potencia menor que $\rho$. Esa frontera ya fue resuelta, porque está
dentro del run $R_{-1}$. Cuando $y$ llegó, $x$ ya estaba en la pila y todas las
entradas sobre $x$ tenían potencia mayor que la de $x$, que a su vez es mayor que la
de $y$. Entonces el lazo de fusiones habría resuelto también a $x$, lo cual es una
contradicción.

(b) Invariante: inmediatamente después de que el run de un elemento $e$ se fusiona
resolviendo una frontera de potencia $q$, toda frontera pendiente adyacente a ese run
tiene potencia menor que $q$:
\begin{itemize}[nosep]
\item la de la izquierda, por (a);
\item la de la derecha, porque es la frontera entrante, que se fusionó justamente por
  ser menor.
\end{itemize}
Las fusiones posteriores conservan esa propiedad: absorber un vecino deja como nueva
frontera adyacente una que es más chica todavía. Así, las potencias resueltas que
involucran a $e$ forman una sucesión estrictamente decreciente de valores en
$[1,P_{\max}]$.

(c) Cada fusión es un SymMerge (Lema~\ref{lem:sym}), y por (b) cada elemento entra en
$O(\log n)$ de ellos. A eso se suman $O(n)$ por detección y $O(32\,n)$ por inserción.
\end{proof}

%======================================================================
\section{Teorema principal}

\begin{teorema}[Corrección total de SotSort]\label{thm:main}
Bajo los Supuestos~\ref{sup:orden}--\ref{sup:n},
\[
\hoare{A=A_0,\ |A_0|=n}{\SotSort(A,n)}{A=\Omega}.
\]
\end{teorema}
\begin{proof}
$A_0$ es coherente porque las etiquetas son los índices. Hay tres casos.
\begin{itemize}[nosep]
\item Si $n<2$ o \Ordenado devuelve verdadero: $A$ queda ordenado y OE
  (Lema~\ref{lem:yaord}), luego coherente. Por el Lema~\ref{lem:coh} con $I=[0,n)$,
  $A=\Omega$.
\item Si \PocosGiros devuelve verdadero: $A$ no cambió desde \Ordenado,
  $A[0..e_0)$ es no decreciente y el Teorema~\ref{thm:natural} da un resultado
  ordenado y OE. Por el Lema~\ref{lem:coh}, $A=\Omega$.
\item En otro caso, $A=A_0$ cumple la precondición del Teorema~\ref{thm:quick} con
  $[0,n)$, que da $A=\Omega$.
\end{itemize}
La terminación se sigue de la de cada componente.
\end{proof}

%======================================================================
\section{Complejidad}

\begin{teorema}[Memoria]\label{thm:mem}
SotSort usa, además del arreglo:
\begin{itemize}[nosep]
\item un buffer de $\Bsz$ elementos;
\item a lo sumo $64+2$ copias de elementos: $\mathit{lb}$ en la pila de \Quick, más
  $p$ y $\mathit{lb}$ actuales;
\item pilas de tamaño acotado: \Quick ($\le 64$), \Part ($\le 57$),
  \SymMerge ($\le 64$) y \Natural ($\le 64$), todas por debajo de las 96 entradas
  reservadas, más $O(1)$ escalares.
\end{itemize}
Todo esto es independiente de $n$ para $n<2^{63}$, y no se usa recursión ni memoria
dinámica.
\end{teorema}
\begin{proof}
Lemas~\ref{lem:partest}, \ref{lem:sym}(a), \ref{lem:prof}(c) y~\ref{lem:natcost}(a).
\end{proof}

\begin{teorema}[Peor caso]\label{thm:cost}
SotSort hace $O(n\log^2 n)$ movimientos y
$O\!\big(n\log n+\frac{n}{\Bsz}\log^3 n\big)$ comparaciones.
\end{teorema}
\begin{proof}
\emph{Camino \Quick.} Por el Lema~\ref{lem:prof}(b), cada elemento pertenece a
$O(\log n)$ intervalos procesados. En cada uno, el intervalo de tamaño $\lambda$
cuesta:
\begin{itemize}[nosep]
\item a lo sumo $2\lambda$ comparaciones de partición (Lema~\ref{lem:partest}) más
  $O(\lambda^{1/3})$ del pivote;
\item $O(\lambda\log(\lambda/\Bsz))$ escrituras.
\end{itemize}
Sumando por niveles de anidamiento, que son disjuntos en cada nivel, se obtienen
$O(n\log n)$ comparaciones y $O(n\log n\log(n/\Bsz))$ escrituras. Las hojas suman
$O(n\log\leaf)$, y los respaldos están acotados por el Lema~\ref{lem:respaldo} sobre
intervalos disjuntos.

\emph{Camino natural.} Lema~\ref{lem:natcost}(c).
\end{proof}

\begin{observacion}
El término $\frac{n}{\Bsz}\log^3 n$ viene de las búsquedas binarias en los nodos
internos de SymMerge. Con $\Bsz=512$ y $n\le 2^{40}$ es menor que $n$, pero como
$\Bsz$ es constante, asintóticamente domina a $n\log n$. No se demuestra aquí si una
variante con la cota $O(m\log(n/m+1))$ de Kim y Kutzner lo elimina.
\end{observacion}

%======================================================================
\section{Modo rápido con memoria adaptable (v8)}\label{sec:rapido}

Las secciones anteriores tratan \SotSort con memoria constante, que en la
implementación de v8 sigue siendo \texttt{stable\_sort}. La versión 8 agrega dos
entradas que usan memoria adicional si la hay: \Rapido (\texttt{stable\_sort\_fast}) y
\PorClave (\texttt{stable\_sort\_by\_key}). Esta sección da su pseudocódigo y las
demostraciones, al mismo nivel de detalle que las anteriores. $\Bsz_0$ denota el buffer
propio de \SotSort (el $\Bsz$ de las secciones anteriores) y $\beta_m$ el presupuesto
de memoria en bytes.

\subsection{Supuestos adicionales y tamaño del buffer}

\begin{supuesto}[Hilos]\label{sup:hilos}
Si $P>1$, $\lt$ puede evaluarse desde varios hilos a la vez, y cada evaluación sigue
cumpliendo el Supuesto~\ref{sup:orden}. Lo que un hilo escribe antes de una barrera (la
espera \texttt{join} de C++) es visible para todos después de ella.
\end{supuesto}
\begin{supuesto}[Enteros]\label{sup:enteros}
Si $T$ es un tipo entero y $\lt$ es $<$, el estado observable de un elemento es su
valor: dos elementos con el mismo valor son indistinguibles. Para \Conteo, ``$A=\Omega$''
significa que $A$ y $\Omega$ tienen la misma secuencia de valores.
\end{supuesto}
\begin{supuesto}[Reservas]\label{sup:reserva}
\Reservar{$x, E$} o falla sin efectos, o entrega $E$ celdas disjuntas de $A$ y de toda
otra memoria en uso. En Linux con \emph{overcommit}, una reserva puede salir bien y
fallar recién al usarse; ese caso queda fuera del modelo.
\end{supuesto}

\begin{observacion}[Buffer más grande que $\Bsz_0$]\label{obs:bsz}
Revisé cada aparición de $\Bsz$ en las demostraciones anteriores. Solo se usan cotas
inferiores ($\Bsz\ge 16$ en la muestra del pivote y en la pila de \Part) y
precondiciones de la forma ``el segmento entra en el buffer'', que con un $\Bsz$ mayor se
cumplen igual o con más razón. La cota superior $\Bsz\le4096$ del
Supuesto~\ref{sup:n} describía la implementación de v7 y no interviene en ninguna
prueba. Por lo tanto los lemas y teoremas anteriores valen para todo $\Bsz\ge 16$, en
particular para $\Bsz=E$. La hoja sigue siendo $\leaf=\min(128,\Bsz_0)\le\Bsz$.
\end{observacion}

% Algoritmos del modo rápido (v8). Se incluye desde sotsort.tex y sotsort_demo.tex.
% Requiere: \lt, \nlt, \Mrg, \buf, \Bsz y las funciones \SotSort, \Ordenado,
% \PocosGiros, \Aleatorio, \MIzq, \MDer, \PlanF, \Rapido, \Galope, \Conteo,
% \CoRango, \Paralelo, \PorClave, \Reservar, \Disp.

\begin{algorithm}[H]
\caption{Modo rápido (v8): plan de memoria y punto de entrada}\label{alg:plan}
\KwIn{$A[0..n)$, comparador $\lt$, opciones $O$: tope $O.\mathit{max}$ en bytes, fracción $O.f$ ($0{,}5$), $O.\mathit{sondear}$, $O.\mathit{hilos}$ ($0$ = automático), $O.\mathit{atajos}$}
\Fn(\tcp*[h]{$s=\mathit{sizeof}(T)$; $\Bsz_0$: buffer propio, unos 16\,KB}){\PlanF{$n, s, O$}}{
  $\beta_m\gets O.\mathit{max}$\;
  \If(\tcp*[f]{solo si el buffer pasaría de 1\,MB}){$ns>2^{20}$ \KwAnd $O.\mathit{sondear}$}{
    $D\gets$ \Disp{}\tcp*{$\min(\text{MemAvailable},\ \text{margen del cgroup})$; $0$ si no se sabe}
    $\beta_m\gets\min\big(\beta_m,\ D>0\ ?\ \lfloor O.f\cdot D\rfloor : 256\,\mathrm{MB}\big)$\;
  }
  $(\mathit{modo},E,P)\gets(\mathsf{O1},0,1)$\;
  \lIf{$n\le\Bsz_0$}{\KwRet $(\mathit{modo},E,P,\beta_m)$}
  \eIf{$ns\le\beta_m$}{$(\mathit{modo},E)\gets(\mathsf{Completo},\,n)$\;}{
    $b\gets\lfloor\beta_m/s\rfloor$;\quad $E'\gets$ mayor potencia de 2 que no supera $b$\;
    \lIf{$E'>\Bsz_0$ \KwAnd $b\ge 2\Bsz_0$}{$(\mathit{modo},E)\gets(\mathsf{Bloques},\,E')$}
  }
  \If{$\mathit{modo}=\mathsf{Completo}$ \KwAnd $O.\mathit{hilos}\ne 1$}{
    $h\gets O.\mathit{hilos}$, o los hilos del hardware si es $0$\;
    $P\gets$ mayor potencia de 2 con $P\le h$ y $P\le\max(1,\lfloor n/2^{17}\rfloor)$\;
  }
  \KwRet $(\mathit{modo},E,P,\beta_m)$\;
}
\BlankLine
\Fn{\Rapido{$A, n, O$}}{
  $\pi\gets$ \PlanF{$n,\ \mathit{sizeof}(T),\ O$}\;
  \lIf{$n<2$}{\KwRet}
  semilla de \Aleatorio${}\gets\sigma\oplus n$, salvo que $O$ pida pivotes aleatorios\;
  \lIf{\Ordenado{$0,n$}}{\KwRet}
  \lIf{$T$ entero, ${\lt}={<}$, $O.\mathit{atajos}$ \KwAnd \Conteo{$A, n, \pi.\beta_m$}}{\KwRet}
  \While(\tcp*[f]{caída escalonada}){$\pi.\mathit{modo}\ne\mathsf{O1}$ \KwAnd \KwNot \Reservar{$\mathit{aux},\ \pi.E$}}{
    $e\gets(\pi.\mathit{modo}=\mathsf{Completo}\ ?\ \lfloor n/2\rfloor : \pi.E/2)$;\quad $E'\gets$ mayor potencia de 2 que no supera $e$;\quad $\pi.P\gets 1$\;
    \leIf{$E'>\Bsz_0$}{$(\pi.\mathit{modo},\pi.E)\gets(\mathsf{Bloques},E')$}{$(\pi.\mathit{modo},\pi.E)\gets(\mathsf{O1},0)$}
  }
  \lIf{$\pi.\mathit{modo}=\mathsf{Completo}$ \KwAnd $\pi.P>1$}{\Paralelo{$A, n, \mathit{aux}, \pi.P$}; \KwRet}
  \lIf{$\pi.\mathit{modo}\ne\mathsf{O1}$}{$\buf\gets\mathit{aux}$;\ $\Bsz\gets\pi.E$;\ $\mathit{galope}\gets\mathbf{true}$}
  \SotSort{$A, n$}\tcp*{el Algoritmo~\ref{alg:main}, con ese $\buf$ y ese $\Bsz$}
}
\end{algorithm}

\noindent $\sigma$ es una constante fija. \Reservar pide memoria sin lanzar
excepciones: devuelve falso si no la obtiene. Los buffers de 4\,MB o más se piden
alineados a 2\,MB y con \texttt{MADV\_HUGEPAGE} (páginas grandes, si el núcleo las da).

\begin{algorithm}[H]
\caption{Mezcla con buffer y galope (v8)}\label{alg:galope}
\Fn(\tcp*[h]{$m-l\le\Bsz$;\ $M_G=7$}){\MIzq{$l, m, r$}}{
  $a\gets m-l$;\quad $\buf[0..a)\gets A[l..m)$;\quad $p\gets 0$;\ $q\gets m$;\ $o\gets l$\;
  \lIf{\KwNot $\mathit{galope}$}{lazo clásico: un paso de cabeza por comparación; \KwRet}
  \While{$p<a$ \KwAnd $q<r$}{
    $c_x\gets 0$;\quad $c_y\gets 0$\tcp*{victorias seguidas de cada lado}
    \While(\tcp*[f]{uno a uno}){$p<a$ \KwAnd $q<r$ \KwAnd $\max(c_x,c_y)<M_G$}{
      \leIf{$A[q]\lt\buf[p]$}{$A[o{+}{+}]\gets A[q{+}{+}]$; $c_y\gets c_y{+}1$; $c_x\gets 0$}{$A[o{+}{+}]\gets\buf[p{+}{+}]$; $c_x\gets c_x{+}1$; $c_y\gets 0$}
    }
    \While(\tcp*[f]{galope}){$p<a$ \KwAnd $q<r$}{
      $k\gets$ \Galope{$\buf[p..a),\ x\mapsto A[q]\nlt x$}\tcp*{lo que no supera a la cabeza de $Y$}
      $A[o..o{+}k)\gets\buf[p..p{+}k)$;\quad $o\gets o{+}k$;\quad $p\gets p{+}k$\;
      \lIf{$p=a$}{\KwBreak}
      $k'\gets$ \Galope{$A[q..r),\ y\mapsto y\lt\buf[p]$}\tcp*{lo estrictamente menor que la cabeza de $X$}
      $A[o..o{+}k')\gets A[q..q{+}k')$;\quad $o\gets o{+}k'$;\quad $q\gets q{+}k'$\tcp*{$o<q$: hacia adelante}
      \lIf{$k<M_G$ \KwAnd $k'<M_G$}{\KwBreak}
    }
  }
  $A[o..o{+}a{-}p)\gets\buf[p..a)$\;
}
\BlankLine
\Fn(\tcp*[h]{$r-m\le\Bsz$}){\MDer{$l, m, r$}}{
  simétrico: la derecha va al buffer, se mezcla desde el final y \Galope busca desde el final\;
}
\BlankLine
\Fn(\tcp*[h]{$Q$ monótono: verdadero y después falso}){\Galope{$S[0..m),\ Q$}}{
  \lIf{$m=0$ \KwOr \KwNot $Q(S[0])$}{\KwRet $0$}
  $i\gets 0$;\quad $\delta\gets 1$\;
  \While(\tcp*[f]{prueba $S[1], S[3], S[7], S[15],\dots$}){$i+\delta<m$ \KwAnd $Q(S[i+\delta])$}{$i\gets i+\delta$;\quad $\delta\gets 2\delta$\;}
  $h\gets\min(i+\delta,\ m)$\;
  \KwRet el menor $j\in(i,h]$ con $j=h$ o $\neg Q(S[j])$, por búsqueda binaria\;
}
\end{algorithm}

\begin{algorithm}[H]
\caption{Conteo de enteros y modo paralelo con \emph{merge path} (v8)}\label{alg:paralelo}
\Fn(\tcp*[h]{$T$ entero y ${\lt}={<}$}){\Conteo{$A, n, \beta_m$}}{
  \lIf{$n<4096$}{\KwRet $\mathbf{false}$}
  \lIf{el rango de $A[0], A[s], A[2s],\dots$ con $s=\lfloor n/1024\rfloor$ es $\ge 2^{16}$}{\KwRet $\mathbf{false}$}
  $\mu\gets\min A$;\quad $R\gets\max A-\mu$\;
  \lIf{$R\ge 2^{16}$ \KwOr $8(R+1)>\beta_m$}{\KwRet $\mathbf{false}$}
  $c[0..R]\gets 0$;\quad \lFor{$i\gets 0$ \KwTo $n-1$}{$c[A[i]-\mu]\gets c[A[i]-\mu]+1$}
  $o\gets 0$;\quad \lFor{$v\gets 0$ \KwTo $R$}{$A[o..o{+}c[v])\gets\mu+v$;\ $o\gets o+c[v]$}
  \KwRet $\mathbf{true}$\;
}
\BlankLine
\Fn(\tcp*[h]{cuántos de $X$ hay entre los primeros $k$ de $\Mrg(X,Y)$}){\CoRango{$k,\ X[0..a),\ Y[0..b)$}}{
  $lo\gets\max(0,k-b)$;\quad $hi\gets\min(k,a)$\;
  \While{$lo<hi$}{
    $i\gets\lfloor(lo+hi)/2\rfloor$;\quad $j\gets k-i$\;
    \leIf(\tcp*[f]{$X[i]$ va antes que $Y[j-1]$}){$Y[j-1]\nlt X[i]$}{$lo\gets i+1$}{$hi\gets i$}
  }
  \KwRet $lo$\;
}
\BlankLine
\Fn(\tcp*[h]{$P$ potencia de 2; $\mathit{aux}$ de $n$ elementos}){\Paralelo{$A, n, \mathit{aux}, P$}}{
  $d_i\gets\lfloor ni/P\rfloor$ para $i=0,\dots,P$\;
  \textbf{en paralelo} para $i<P$: \SotSort sobre $A[d_i..d_{i+1})$ con $\buf=\mathit{aux}[d_i..d_{i+1})$ y galope\;
  $S\gets A$;\quad $D\gets\mathit{aux}$;\quad $w\gets 1$\;
  \For{$\rho\gets P, P/2,\dots,2$}{
    $g\gets 2P/\rho$\tcp*{tareas por cada par de runs}
    \ForEach(\tcp*[f]{par $u=\lfloor t/g\rfloor$, parte $v=t\bmod g$}){tarea $t<P$}{
      $lo\gets d_{2uw}$;\ $md\gets d_{(2u+1)w}$;\ $hi\gets d_{(2u+2)w}$;\ $N\gets hi-lo$\;
      $k_0\gets\lfloor Nv/g\rfloor$;\ $k_1\gets\lfloor N(v{+}1)/g\rfloor$;\ $i_\ell\gets$ \CoRango{$k_\ell,\ S[lo..md),\ S[md..hi)$} para $\ell=0,1$\;
    }
    \textbf{en paralelo} para cada tarea: mezclar $S[lo{+}i_0..lo{+}i_1)$ con $S[md{+}k_0{-}i_0..md{+}k_1{-}i_1)$ en $D[lo{+}k_0..lo{+}k_1)$\;
    intercambiar $S$ y $D$;\quad $w\gets 2w$\;
  }
  \lIf{$S\ne A$}{\textbf{en paralelo}, copiar $S$ a $A$ por trozos}
}
\end{algorithm}

\noindent Los cortes de cada ronda se calculan todos antes de mezclar: con tipos que
no se copian byte a byte, mover vacía el origen, y otro hilo podría estar buscando su
corte sobre esos mismos elementos (ThreadSanitizer encontró esa carrera en una
versión anterior del código).

\begin{algorithm}[H]
\caption{Orden estable por clave entera (v8)}\label{alg:clave}
\KwIn{$A[0..n)$ y una clave entera $\kappa$; el orden es $x\lt_\kappa y\iff\kappa(x)<\kappa(y)$}
\Fn{\PorClave{$A, n, \kappa, O$}}{
  $\pi\gets$ \PlanF{$n,\ \mathit{sizeof}(T),\ O$}\;
  \If{$T$ no se copia byte a byte \KwOr $n<4096$ \KwOr $\pi.\mathit{modo}\ne\mathsf{Completo}$}{\Rapido{$A, n, O$} con $\lt_\kappa$;\quad \KwRet\;}
  \lIf{$\kappa(A[0]),\dots,\kappa(A[n-1])$ es no decreciente}{\KwRet}
  \lIf{\PocosGiros{$n$}}{\Rapido{$A, n, O$} con $\lt_\kappa$; \KwRet}
  $\mu\gets\min_i\kappa(A[i])$;\quad $R\gets\max_i\kappa(A[i])-\mu$;\quad $r\gets\mathrm{bits}(R)$;\quad $\iota\gets\mathrm{bits}(n-1)$\;
  \If(\tcp*[f]{indirecto empaquetado}){$\mathit{sizeof}(T)\ge 32$ \KwAnd $r+\iota\le 64$ \KwAnd $n(\mathit{sizeof}(T)+16)\le\pi.\beta_m$}{
    $K[i]\gets(\kappa(A[i])-\mu)\cdot 2^{\iota}+i$ para $i<n$\tcp*{enteros de 64 bits distintos}
    \Rapido{$K, n$} con $<$, sin sondeo y con tope de $8n$ bytes\;
    $\mathit{aux}[j]\gets A[K[j]\bmod 2^{\iota}]$ para $j<n$;\quad $A\gets\mathit{aux}$;\quad \KwRet\;
  }
  \If(\tcp*[f]{radix LSD sobre los registros}){$R<2^{24}$}{
    $q\gets 1$, $2$ o $3$ según $R<2^8$, $2^{16}$ o $2^{24}$;\quad $S\gets A$;\quad $D\gets\mathit{aux}$\;
    \For{$\delta\gets 0$ \KwTo $q-1$}{distribuir $S$ en $D$, en orden, por el byte $\delta$ de $\kappa(x)-\mu$;\quad intercambiar $S$ y $D$\;}
    \lIf{$S\ne A$}{$A\gets S$}
    \KwRet\;
  }
  \Rapido{$A, n, O$} con $\lt_\kappa$\;
}
\end{algorithm}

\noindent $\mathrm{bits}(x)$ es la cantidad de bits de $x$ ($\mathrm{bits}(0)=0$). Si una
reserva falla, se pasa a la rama siguiente.


\subsection{Mezcla con galope}

\begin{lema}[Galope]\label{lem:galope}
Sea $Q$ monótono sobre $S[0..m)$: existe $k^*$ con $Q(S[i])\iff i<k^*$. Entonces
\Galope{$S,Q$} devuelve $k^*$ y evalúa $Q$ a lo sumo $2\lfloor\log_2 k^*\rfloor+2$ veces
si $k^*\ge1$, y una vez si $k^*=0$ (ninguna si $m=0$). En los dos casos son a lo sumo
$2\lfloor\log_2(k^*{+}1)\rfloor+2$.
\end{lema}
\begin{proof}
Si $m=0$ o $\neg Q(S[0])$, $k^*=0$. Si no, el lazo exponencial mantiene $Q(S[i])$, es
decir, $i<k^*$; la variante es $m-i$, porque $i$ crece en cada vuelta. Al salir, o bien
$i+\delta\ge m$ y entonces $k^*\le m=h$, o bien $\neg Q(S[i+\delta])$ y entonces
$k^*\le i+\delta=h$. En los dos casos $k^*\in(i,h]$. Para $j\in(i,k^*)$ vale $Q(S[j])$, y
en $j=k^*$ vale $j=h$ o $\neg Q(S[j])$: el menor $j$ de la línea final es $k^*$, y la
búsqueda binaria lo encuentra porque la condición es monótona en $j$.

\emph{Costo.} Las posiciones probadas por el lazo son $2^{t+1}-1$ para $t=0,1,\dots$;
las que dan verdadero son las menores que $k^*$, o sea las $t$ con
$2^{t+1}\le k^*$: a lo sumo $\lfloor\log_2k^*\rfloor$. Hay a lo sumo una prueba que da
falso. El intervalo que queda para la búsqueda binaria tiene a lo sumo
$\delta-1=2^{\lfloor\log_2 k^*\rfloor}-1$ posiciones, que se resuelven con a lo sumo
$\lfloor\log_2k^*\rfloor$ evaluaciones. Con la evaluación de $S[0]$, el total es a lo
sumo $2\lfloor\log_2k^*\rfloor+2$. La búsqueda desde el final de \MDer es simétrica,
con posiciones $1, 2, 4, 8,\dots$ contadas desde el final, y tiene la misma cota.
\end{proof}

\begin{lema}[Monotonía]\label{lem:galmono}
Si $X$ está ordenada, $x\mapsto y\nlt x$ es monótono sobre $X$ para todo $y$. Si $Y$
está ordenada, $y\mapsto y\lt x$ es monótono sobre $Y$ para todo $x$.
\end{lema}
\begin{proof}
Sean $j<i$. Como $X$ está ordenada, $X[j]\les X[i]$. Si $y\nlt X[i]$, es decir,
$X[i]\les y$, la transitividad de $\les$ (Lema~\ref{lem:swo}(d)) da $X[j]\les y$. En
el segundo caso, si $Y[i]\lt x$, el Lema~\ref{lem:swo}(c) con $Y[j]\les Y[i]$ da
$Y[j]\lt x$.
\end{proof}

\begin{lema}[MergeIzq con galope]\label{lem:mizqg}
Con $\mathit{galope}$ verdadero, \MIzq cumple la misma tripla del
Lema~\ref{lem:mizq}, con a lo sumo $6(|X|+|Y|)$ comparaciones. Lo mismo vale para
\MDer.
\end{lema}
\begin{proof}
Se usa el invariante del Lema~\ref{lem:mizq}: $o=l+p+(q-m)$,
$A[l..o)=\Mrg(X[0..p),Y[0..q{-}m))$, $\buf[p..a)=X[p..)$ y $A[q..r)=Y[q{-}m..)$. Como
$o-q=p-a$, vale $o\le q$, con desigualdad estricta mientras $p<a$.

Cada paso del lazo uno a uno es un paso de cabeza (Lema~\ref{lem:merge}(a)). En el
galope sobre $X$, la cabeza de $Y$ es fija, $y=A[q]$, y por los
Lemas~\ref{lem:galmono} y~\ref{lem:galope} los $k$ elementos copiados son exactamente
los primeros de $X[p..)$ que cumplen $y\nlt x$. Para cada uno, el paso de cabeza elige
$X$, porque no vale $Y[\cdot]\lt X[\cdot]$; así, la copia equivale a $k$ pasos de
cabeza. Si después $p<a$, el elemento $X[p]$ no cumple el predicado: $A[q]\lt X[p]$. En
el galope sobre $Y$ la cabeza de $X$ es fija, $x=\buf[p]$, y los $k'$ elementos
copiados son los primeros de $Y[q{-}m..)$ estrictamente menores que $x$: son $k'$ pasos
de cabeza que eligen $Y$. Como $A[q]\lt\buf[p]$, vale $k'\ge 1$. La copia de $Y$ va
hacia adelante con $o<q$, así que no pisa nada sin leer.

\emph{Terminación.} Cada vuelta del lazo externo hace al menos un paso uno a uno, y
cada vuelta del galope o deja $p=a$ o copia $k'\ge1$ elementos: la variante
$(a-p)+(r-q)$ decrece. Al salir, el final es el del Lema~\ref{lem:mizq}.

\emph{Comparaciones.} En el lazo uno a uno, una por elemento escrito. Una vuelta del
galope que hace las dos búsquedas cuesta, por el Lema~\ref{lem:galope}, a lo sumo
$2\lfloor\log_2(k{+}1)\rfloor+2+2\lfloor\log_2(k'{+}1)\rfloor+2\le 2k+2k'+4\le 6(k+k')$,
porque $k'\ge1$. Una vuelta que termina con $p=a$ copió $k\ge1$ elementos y cuesta a lo
sumo $2k+2\le 4k$. Cada elemento se escribe una sola vez, así que el total es a lo sumo
$6(|X|+|Y|)$. \MDer es simétrico, con los pasos de cola del Lema~\ref{lem:merge}(b).
\end{proof}

\begin{observacion}
Con galope, en el Lema~\ref{lem:sym}(d) las mezclas de las hojas comparan a lo sumo 6
veces por elemento en lugar de una: la cota pasa a
$6s+\bigl(1+\frac{2s}{\Bsz}\bigr)(5\ell+2)$. Las cotas asintóticas del
Lema~\ref{lem:respaldo}, del Lema~\ref{lem:natcost} y del Teorema~\ref{thm:cost} no
cambian; solo sus constantes. El galope está apagado en el modo $O(1)$ para tipos
aritméticos, que es como se midió y demostró v7.
\end{observacion}

\subsection{Conteo de enteros}

\begin{lema}[Conteo]\label{lem:conteo}
Bajo el Supuesto~\ref{sup:enteros},
\[
\hoare{A=A_0}{c\gets\Conteo(A,n,\beta_m)}{(\neg c\wedge A=A_0)\ \vee\ (c\wedge A=\Omega)},
\]
con $O(n+R)$ operaciones y $8(R+1)\le\beta_m$ bytes de contadores, que se liberan antes
de volver.
\end{lema}
\begin{proof}
Hasta la línea 5, \Conteo solo lee. La muestra es un subconjunto de $A$, así que su
rango no supera a $R$: descartar por la muestra es descartar un caso con $R\ge2^{16}$.
Si devuelve falso, $A=A_0$.

Las diferencias $A[i]-\mu$ se calculan en aritmética sin signo de $w$ bits, módulo
$2^w$. Como $0\le A[i]-\mu\le R\le 2^w-1$, el resultado es la diferencia verdadera y el
índice está en $[0,R]$.

Invariante del primer lazo: tras $i$ vueltas, $c[v]=\#\{j<i:\ A_0[j]=\mu+v\}$. Invariante
del segundo: tras procesar los valores $0,\dots,v-1$, vale $o=\sum_{u<v}c[u]$ y $A[0..o)$
es la secuencia no decreciente que repite cada $\mu+u$, para $u<v$, exactamente $c[u]$
veces. Al final $o=n$ y $A$ es la única secuencia no decreciente con el multiconjunto de
valores de $A_0$. $\Omega$ está ordenada, que para enteros con $<$ significa no
decreciente, y tiene ese mismo multiconjunto: las dos secuencias de valores coinciden.
Por el Supuesto~\ref{sup:enteros}, $A=\Omega$.
\end{proof}

\subsection{Modo paralelo}

\begin{lema}[CoRango]\label{lem:corango}
Sean $X[0..a)$ e $Y[0..b)$ ordenadas, $Z=\Mrg(X,Y)$, $0\le k\le a+b$, y sea $c(k)$ la
cantidad de elementos de $X$ en $Z[0..k)$. Entonces \CoRango{$k,X,Y$} devuelve $c(k)$
con a lo sumo $\lceil\log_2(\min(k,a)-\max(0,k-b)+1)\rceil$ comparaciones, $c$ es no
decreciente, y para $k_0\le k_1$:
\[
Z[k_0..k_1)=\Mrg\big(X[c(k_0)..c(k_1)),\ Y[k_0{-}c(k_0)..k_1{-}c(k_1))\big).
\]
\end{lema}
\begin{proof}
Como $\Mrg$ preserva el orden interno de $X$ y de $Y$, los elementos de $X$ en
$Z[0..k)$ son $X[0..c)$ y los de $Y$ son $Y[0..k{-}c)$, con $c=c(k)$. Por eso
$\max(0,k-b)\le c\le\min(k,a)$.

\emph{Afirmación.} Para $i\in[\max(0,k-b),\min(k,a))$ y $j=k-i$ (vale $1\le j\le b$):
$i<c\iff Y[j-1]\nlt X[i]$.
($\Rightarrow$) Si $i<c$, $X[i]$ está en $Z[0..k)$. Como $k-c<k-i=j$, $Y[j-1]$ no está
en $Z[0..k)$, así que $X[i]$ precede a $Y[j-1]$ en $Z$, que está ordenada:
$Y[j-1]\nlt X[i]$.
($\Leftarrow$) Si $i\ge c$, $X[i]$ no está en $Z[0..k)$ y, como $k-c\ge j$, $Y[j-1]$
sí. Entonces $Y[j-1]$ precede a $X[i]$ y, por estar ordenada, $Y[j-1]\les X[i]$. Si
fueran equivalentes, la Definición~\ref{def:merge} pondría a $X[i]$ primero; luego
$Y[j-1]\lt X[i]$.

La búsqueda binaria mantiene $lo\le c\le hi$: si la condición vale, $i<c$ y
$lo\gets i+1\le c$; si no, $i\ge c$ y $hi\gets i$. Los índices son válidos porque
$lo\le i<hi$ da $0\le i<a$ y $1\le j\le b$. La variante es $hi-lo$, que al menos se
reduce a la mitad, lo que da la cota de comparaciones. Al terminar, $lo=hi=c$.

$c$ es no decreciente porque $Z[0..k)$ crece con $k$. La igualdad final es el
Lema~\ref{lem:merge}(c): $Z[k_0..k_1)$ es un segmento de $Z$ y contiene exactamente esos
elementos de $X$ y de $Y$.
\end{proof}

\begin{lema}[Paralelo]\label{lem:paralelo}
Bajo los Supuestos~\ref{sup:hilos} y~\ref{sup:reserva}, con $P=2^p\le n$ y
$\mathit{aux}$ de $n$ celdas:
\[
\hoare{A=A_0}{\Paralelo(A,n,\mathit{aux},P)}{A=\Omega}.
\]
Además, ninguna celda es escrita por dos hilos en la misma fase, y ninguna celda que un
hilo lee en una fase es escrita por otro en esa fase.
\end{lema}
\begin{proof}
\emph{Fase de trozos.} El hilo $i$ solo toca $A[d_i..d_{i+1})$ y
$\mathit{aux}[d_i..d_{i+1})$, que son disjuntos entre hilos. Las llamadas a $\lt$ son
seguras por el Supuesto~\ref{sup:hilos}. Por el Teorema~\ref{thm:main}, con
$\Bsz=d_{i+1}-d_i$ (Observación~\ref{obs:bsz}) o con el buffer propio si el trozo es más
chico, cada trozo queda ordenado y la transformación es OE dentro del trozo. Por el
Lema~\ref{lem:oe}(c), es OE sobre todo el arreglo, que queda coherente.

\emph{Invariante de rondas.} Al empezar la ronda de ancho $w$, con $\rho=P/w$ runs:
$S$ contiene cada elemento de $A_0$ una vez, es coherente, y cada run
$S[d_{jw}..d_{(j+1)w})$ está ordenado. Sea un par $X=S[lo..md)$, $Y=S[md..hi)$ y
$Z=\Mrg(X,Y)$. Sus $g$ tareas usan los cortes $\lfloor Nv/g\rfloor$, $v=0,\dots,g$, que
parten $[0,N)$ en tramos consecutivos. Por el Lema~\ref{lem:corango}, la tarea $v$
escribe $D[lo{+}k_0..lo{+}k_1)=Z[k_0..k_1)$, y entre todas escriben
$D[lo..hi)=Z$. $\Mrg$ es OE (Lema~\ref{lem:merge}(d)) y los pares son disjuntos, así que
$D$ es coherente y sus runs de ancho $2w$ están ordenados. Al intercambiar, el invariante
vale para $2w$.

\emph{Carreras.} En el cálculo de cortes solo se lee $S$. En la mezcla, cada tarea
escribe un tramo de $D$ disjunto de los demás y solo lee $S$, que nadie escribe en esa
fase: $A$ y $\mathit{aux}$ son disjuntos (Supuesto~\ref{sup:reserva}). Las barreras
separan las fases (Supuesto~\ref{sup:hilos}). La copia final escribe tramos disjuntos de
$A$ y lee solo $S$. Que los cortes se calculen todos antes de mezclar es necesario:
para tipos no triviales, mover deja el origen en un estado no especificado, y una
búsqueda de corte concurrente leería esos valores.

\emph{Final.} Tras $p$ rondas queda un solo run: $S[0..n)$ está ordenado y es coherente,
así que $S=\Omega$ por el Lema~\ref{lem:coh}. Si $S\ne A$, la copia lo lleva a $A$
preservando identidades (Supuesto~\ref{sup:ident}). Termina porque hay $p$ rondas y cada
lazo es finito.
\end{proof}

\subsection{Corrección, memoria y cotas del modo rápido}

\begin{teorema}[Corrección de \Rapido]\label{thm:rapido}
Bajo los Supuestos~\ref{sup:orden}--\ref{sup:n} y~\ref{sup:reserva}, más el
Supuesto~\ref{sup:hilos} si $P>1$ y el Supuesto~\ref{sup:enteros} si se usa \Conteo:
\[
\hoare{A=A_0,\ |A_0|=n}{\Rapido(A,n,O)}{A=\Omega}.
\]
\end{teorema}
\begin{proof}
\PlanF solo calcula números y no toca $A$. Si $n<2$ no hay nada que hacer. Si \Ordenado
devuelve verdadero, $A$ queda ordenado y OE (Lema~\ref{lem:yaord}) y, por el
Lema~\ref{lem:coh}, $A=\Omega$. Si \Conteo devuelve verdadero, $A=\Omega$ por el
Lema~\ref{lem:conteo}; si devuelve falso, $A=A_0$. El lazo de reservas termina porque
el modo baja estrictamente ($\mathsf{Completo}\to\mathsf{Bloques}$ con $E$ menor o
$\mathsf{O1}$, y en $\mathsf{Bloques}$ $E$ se reduce a la mitad) y en $\mathsf{O1}$ se
detiene; no modifica $A$. Con $P>1$, $P\le n/2^{17}$ por \PlanF y el
Lema~\ref{lem:paralelo} da $A=\Omega$. En otro caso se ejecuta \SotSort con $\Bsz=E$ o
$\Bsz=\Bsz_0$, y con galope si hay buffer: el Teorema~\ref{thm:main} vale por la
Observación~\ref{obs:bsz} y el Lema~\ref{lem:mizqg}. \SotSort vuelve a llamar a
\Ordenado, que ya devolvió falso sobre el mismo $A$ y vuelve a hacerlo sin modificarlo
(Lema~\ref{lem:yaord}).
\end{proof}

\begin{teorema}[Memoria de \Rapido]\label{thm:rapmem}
Además del arreglo, \Rapido usa:
\begin{enumerate}[label=(\alph*),nosep]
\item a lo sumo $\beta_m$ bytes reservados en todo momento: el buffer
  $\mathit{aux}$, de $E\cdot\mathit{sizeof}(T)\le\beta_m$ bytes, o los contadores de
  \Conteo, de $8(R+1)\le\beta_m$ bytes, nunca los dos a la vez;
\item la memoria fija del Teorema~\ref{thm:mem} (buffer propio y pilas) por cada hilo
  que ordena un trozo;
\item $O(P)$ palabras para los bordes, los cortes y los hilos, más la pila que el
  sistema le asigna a cada uno de los $P-1$ hilos adicionales.
\end{enumerate}
\end{teorema}
\begin{proof}
(a) En modo completo, $ns\le\beta_m$. En modo bloques, $E\le b=\lfloor\beta_m/s\rfloor$,
así que $Es\le\beta_m$. En la caída escalonada, desde completo
$E'\le\lfloor n/2\rfloor$ y desde bloques $E'\le E/2$: el tamaño solo baja. \Conteo
verifica $8(R+1)\le\beta_m$ y libera los contadores antes de volver, y la reserva del
buffer ocurre después. (b) y (c) son inspección del Algoritmo~\ref{alg:paralelo}.
\end{proof}

\noindent El presupuesto es $\beta_m=\min(O.\mathit{max},\ O.f\cdot D)$, con $D$ la
memoria disponible \emph{en el momento de la consulta}. El teorema garantiza que no se
pide más que eso; no garantiza que otros procesos no consuman la memoria restante
mientras se ordena.

\begin{teorema}[Cotas con buffer completo]\label{thm:rapcost}
Si \SotSort se ejecuta con $\Bsz\ge n$, hace $O(n\log n)$ comparaciones y $O(n\log n)$
movimientos en el peor caso, con galope o sin él. \Paralelo con $\mathit{aux}$ de $n$
celdas hace en total $O(n\log n)$ comparaciones y movimientos.
\end{teorema}
\begin{proof}
\emph{Camino \Quick.} Por el Lema~\ref{lem:prof}(b), los intervalos procesados que
contienen a un elemento forman una cadena de longitud $O(\log n)$. En un intervalo de
tamaño $\lambda\le n\le\Bsz$, cada \Part tiene $k=\lceil\lambda/\Bsz\rceil=1$ bloque: por
el Lema~\ref{lem:partest}, $\lambda$ comparaciones, a lo sumo $3\lambda$ escrituras y
ninguna rotación. Cada iteración hace a lo sumo dos particiones y elige un pivote con
$O(\lambda^{1/3}\log\lambda)$ operaciones. Los intervalos de una misma profundidad son
disjuntos, así que la suma es $O(n\log n)$. Las hojas son disjuntas y cuestan
$O(\leaf\log\leaf)$ cada una (Lema~\ref{lem:hoja}).

\emph{Respaldo y SymMerge.} En todo SymMerge de tamaño $s\le n\le\Bsz$ vale
$\min(n_l,n_r)\le\Bsz$, así que la raíz es una hoja y no hay división: a lo sumo $3s$
escrituras (Lema~\ref{lem:sym}(c)) y $6s+2\lceil\log_2(s{+}1)\rceil+1$ comparaciones
(Lema~\ref{lem:mizqg} y las búsquedas del recorte). \Respaldo sobre un intervalo de
tamaño $\lambda$ hace $O(\lambda)$ por nivel y $O(\log\lambda)$ niveles, más la inserción
en bloques de 24; los intervalos de respaldo son disjuntos.

\emph{Camino natural.} Por el Lema~\ref{lem:natcost}(b), cada elemento participa en
$O(\log n)$ fusiones, cada una lineal por lo anterior; la detección y la inserción de
runs cortos son $O(n)$.

\emph{Paralelo.} Los trozos cuestan $O(\frac{n}{P}\log\frac{n}{P})$ cada uno por lo
anterior. Cada una de las $\log_2P$ rondas hace $n$ movimientos, a lo sumo $n$
comparaciones de mezcla y $2P$ llamadas a \CoRango de $O(\log n)$ comparaciones. Como
$P\le n$, el total es $O(n\log n)$.
\end{proof}

\noindent En modo bloques valen las cotas del Teorema~\ref{thm:cost} con $\Bsz=E$:
$O(n\log n\log(n/E))$ movimientos y $O(n\log n+\frac{n}{E}\log^2n)$ comparaciones.

\begin{teorema}[Determinismo]\label{thm:det}
\begin{enumerate}[label=(\alph*),nosep]
\item El resultado es $\Omega$ cualquiera sea el plan, la semilla, la cantidad de hilos
  y el resultado de las reservas.
\item Si $O.\mathit{sondear}$ es falso, el plan depende solo de $n$,
  $\mathit{sizeof}(T)$, $\Bsz_0$ y $O$, y además de la cantidad de hilos del hardware si
  $O.\mathit{hilos}=0$.
\item Fijado el plan, si ninguna reserva falla y los pivotes no son aleatorios, la
  secuencia de operaciones de cada hilo es función de la entrada. Solo varía el
  entrelazado en el tiempo entre hilos, que no cambia ningún valor leído.
\end{enumerate}
\end{teorema}
\begin{proof}
(a) Es el Teorema~\ref{thm:rapido}. (b) \PlanF solo llama a \Disp si
$O.\mathit{sondear}$ es verdadero, y solo consulta el hardware si $O.\mathit{hilos}=0$.
(c) Cada decisión depende de $A$, del plan y del generador, cuya semilla es
$\sigma\oplus n$ (y una semilla fija por trozo). En \Paralelo, por el
Lema~\ref{lem:paralelo}, cada tarea lee solo valores que no cambian durante su fase.
\end{proof}

\subsection{Orden por clave}

Sea $\kappa$ una clave entera y $x\lt_\kappa y\iff\kappa(x)<\kappa(y)$, que es un orden
débil estricto porque $<$ lo es sobre los enteros. $\Omega$ se toma respecto de
$\lt_\kappa$, con etiquetas iguales a los índices.

\begin{lema}[Indirecto empaquetado]\label{lem:indirecto}
Sean $d_i=\kappa(A_0[i])-\mu\in[0,R]$, $r=\mathrm{bits}(R)$, $\iota=\mathrm{bits}(n-1)$
con $r+\iota\le 64$, y $K[i]=d_i\cdot2^\iota+i$. Entonces los $K[i]$ son distintos,
menores que $2^{64}$, y $K[i]<K[j]\iff A_0[i]\lts A_0[j]$. En consecuencia, si $K$ queda
ordenado de forma creciente, $A[j]\gets A_0[K[j]\bmod 2^\iota]$ deja $A=\Omega$.
\end{lema}
\begin{proof}
$i<2^\iota$ y $d_i<2^r$, así que $K[i]<2^{r+\iota}\le2^{64}$, y $K[i]\bmod2^\iota=i$,
$\lfloor K[i]/2^\iota\rfloor=d_i$: la codificación es inyectiva. Como $i,j<2^\iota$,
comparar $K$ es comparar lexicográficamente $(d,\text{índice})$:
$K[i]<K[j]\iff d_i<d_j\vee(d_i=d_j\wedge i<j)$. Además $d_i<d_j\iff A_0[i]\lt_\kappa
A_0[j]$ y $d_i=d_j\iff A_0[i]\eqv A_0[j]$, que con etiquetas iguales a los índices es la
Definición~\ref{def:omega} de $\lts$. Las restas se hacen sin signo, módulo $2^w$, y son
exactas porque $0\le\kappa-\mu\le R<2^w$. $K$ tiene valores distintos, así que su orden
creciente es único; el recorrido final pone los elementos en orden creciente según
$\lts$, que es $\Omega$. $\mathit{aux}$ es disjunto de $A$ y la copia preserva
identidades.
\end{proof}

\begin{lema}[Radix por clave]\label{lem:radix}
Sea $q$ tal que $R<2^{8q}$ y $\kappa_\delta(x)=(\kappa(x)-\mu)\bmod 2^{8(\delta+1)}$.
Tras la pasada $\delta$ (contando desde 0), $S$ contiene cada elemento una vez, en orden
creciente según el par $(\kappa_\delta(x),\etq(x))$. Al final, $A=\Omega$.
\end{lema}
\begin{proof}
Antes de la primera pasada, $S=A_0$ está en orden creciente de etiqueta. Una pasada
reparte por el byte $\delta$ con sumas de prefijos y un recorrido en orden de $S$: los
elementos quedan en orden creciente de byte y, a igual byte, en el orden que tenían en
$S$. Como $\kappa_\delta=\text{byte}_\delta\cdot2^{8\delta}+\kappa_{\delta-1}$, el
nuevo orden es creciente según $(\text{byte}_\delta,\kappa_{\delta-1},\etq)$, es decir,
según $(\kappa_\delta,\etq)$. Tras $q$ pasadas, $\kappa_{q-1}=\kappa-\mu$ porque
$R<2^{8q}$, y ese orden es $\lts$. Si el resultado quedó en $\mathit{aux}$, se copia.
\end{proof}

\begin{teorema}[\PorClave]\label{thm:clave}
Bajo los mismos supuestos que el Teorema~\ref{thm:rapido},
\[\hoare{A=A_0}{\PorClave(A,n,\kappa,O)}{A=\Omega},\]
y la memoria reservada en todo
momento es a lo sumo $\beta_m$ bytes más $O(1)$ (los contadores del radix).
\end{teorema}
\begin{proof}
Las ramas que delegan en \Rapido con $\lt_\kappa$ son correctas por el
Teorema~\ref{thm:rapido}. Si las claves ya son no decrecientes, $A_0$ está ordenada y es
coherente: $A_0=\Omega$. La rama indirecta es correcta por el Lema~\ref{lem:indirecto}
y el Teorema~\ref{thm:rapido} aplicado a $K$; la de radix, por el
Lema~\ref{lem:radix}. Memoria: la rama indirecta tiene a la vez $K$ ($8n$ bytes),
$\mathit{aux}$ ($n\,\mathit{sizeof}(T)$) y lo que use \Rapido sobre $K$, acotado por su
tope de $8n$ (Teorema~\ref{thm:rapmem}); en total $n(\mathit{sizeof}(T)+16)\le\beta_m$,
que la rama verifica. La de radix usa $\mathit{aux}$, que entra porque el modo es
completo, y $3\cdot256$ contadores.
\end{proof}


%======================================================================
\section{Versión 7 y reconstrucción v8-r}\label{sec:v8r}

\paragraph{Estado.} El código fuente de la versión 8 se perdió. La implementación
actual (\texttt{sotsort.hpp}, ``v8-r'') se reescribió a partir de este documento, del
libro y de una versión anterior del código. Las demostraciones de las secciones
anteriores son sobre el pseudocódigo y no cambian; lo que cambia es que la
correspondencia con el código vuelve a apoyarse solo en pruebas diferenciales contra
\texttt{std::stable\_sort} (sección~\ref{sec:v8r-pruebas}). Esta sección agrega lo que
el documento no cubría: dos cambios de la versión 7 y cuatro adiciones de v8-r. Las
cuatro adiciones solo eligen entre procedimientos ya demostrados o reordenan
recorridos; ninguna introduce un procedimiento de ordenamiento nuevo.

\subsection{Versión 7: pivote con posiciones pseudoaleatorias}
\Pivote toma sus muestras en posiciones $l+js+\Aleatorio{}\bmod s$ en lugar de $l+js$,
con un generador determinista sembrado con $\sigma\oplus n$. El Lema~\ref{lem:pivote}
solo usa que $p$ es copia de un elemento de \seg{A}{l}{r}, y cada posición muestreada
cumple $l\le l+js+(\Aleatorio{}\bmod s)<r$ porque $(j+1)s\le r-l$. El lema y todo lo que
depende de él valen sin cambios. El Teorema~\ref{thm:det} tampoco cambia: la secuencia de
\Aleatorio es función de $n$ y de la semilla.

\subsection{Versión 7: partición de bloque vectorial}

\begin{algorithm}[H]
\caption{Partición de bloque vectorial, con $w$ carriles}\label{alg:partvec}
\KwIn{$s<e$ con $L=e-s\le\Bsz$; pivote $p$; $\triangleleft\in\{\lt,\les\}$; $T$ entero y $\lt$ es $<$}
\Fn{\PartVec{$s, e, p, \triangleleft$}}{
  $m\gets 0$;\quad $i\gets 0$\;
  \While{$i+w\le L$}{
    $X\gets A[s+i\,..\,s+i+w)$\tcp*{una carga}
    $k\gets\{\,j<w : X[j]\triangleleft p\,\}$\;
    $A[s+m\,..\,s+m+w)\gets\mathrm{comprimir}(X,k)$\;
    $\buf[i-m\,..\,i-m+w)\gets\mathrm{comprimir}(X,\overline{k})$\;
    $m\gets m+|k|$;\quad $i\gets i+w$\;
  }
  cola escalar: el lazo de \PartBloque desde $i$\;
  $A[s+m\,..\,e)\gets\buf[0\,..\,L-m)$;\quad \KwRet $s+m$\;
}
\end{algorithm}

\noindent $\mathrm{comprimir}(X,k)$ es el vector con los carriles de $X$ marcados en $k$,
en su orden, seguidos de $w-|k|$ carriles de contenido arbitrario.

\begin{lema}[Partición vectorial]\label{lem:partvec}
Con $C=\seg{A_0}{s}{e}$ el contenido inicial y $P(x)\equiv x\triangleleft p$, al terminar
$\seg{A}{s}{e}=\filt(P,C)\cat\filt(\neg P,C)$ y se devuelve $s+|\filt(P,C)|$. Ninguna
escritura cae fuera de \seg{A}{s}{e} ni de $\buf[0..L)$.
\end{lema}
\begin{proof}
Invariante del lazo: $\seg{A}{s}{s+m}=\filt(P,C[0..i))$,
$\buf[0..i-m)=\filt(\neg P,C[0..i))$ y $\seg{A}{s+i}{e}=C[i..L)$.
\emph{Preservación.} $X=C[i..i+w)$ por la tercera cláusula. El store sobre $A$ escribe
$[s+m,s+m+w)$; como $m\le i$, ese rango está contenido en $[s,s+i+w)$, cuyo contenido
ya fue leído ($C[0..i)$ en vueltas anteriores, $C[i..i+w)$ en $X$), así que no destruye
nada de $C[i+w..L)$, y $s+m+w\le s+i+w\le e$. Sus primeros $|k|$ carriles son
$\filt(P,X)$ y extienden el prefijo. El store sobre $\buf$ escribe
$[i-m,i-m+w)$ con $i-m+w\le i+w\le L\le\Bsz$; sus primeros $w-|k|$ carriles son
$\filt(\neg P,X)$. \emph{Terminación.} La cola escalar es el lazo de
\PartBloque con el mismo invariante (Lema~\ref{lem:partbloque}), y la copia final es
la suya.
\end{proof}

Por el Supuesto~\ref{sup:enteros}, evaluar $\triangleleft$ con una comparación vectorial
de enteros da el mismo resultado que $\lt$ carril por carril. El
Lema~\ref{lem:partest} usa a \PartBloque solo a través de su poscondición, que
\PartVec comparte.

\subsection{v8-r: sonda de rachas en SymMerge}\label{sec:sonda}
En una hoja de \SymMerge (un lado entra en el buffer) hay tres mezclas posibles: copiar
todo al buffer y hacer un \Parity de largos distintos, si los dos runs entran juntos;
copiar el lado corto y usar \MIzq/\MDer con galope; o \MIzq/\MDer con el lazo
clásico. La primera no tiene saltos y conviene cuando los runs se intercalan elemento
por elemento; el galope conviene con rachas largas. En el camino de runs naturales,
para tipos que se copian byte a byte y $a+b\ge256$, se decide con una sonda.

Sean $X=\seg{A}{l}{m}$ e $Y=\seg{A}{m}{r}$ ordenados, $a=|X|$, $b=|Y|$. La sonda toma
$x_0=l$, $x_i=l+\lfloor ia/9\rfloor$ para $1\le i\le 8$, $x_9=m$, y
$y_0=m$, $y_i=\min\{\,j\in[y_{i-1},r] : j=r \vee A[j]\nlt A[x_i]\,\}$ para $1\le i\le8$,
$y_9=r$. Define
\[\mathit{sw}=\sum_{i=1}^{9}\bigl(2\min(x_i-x_{i-1},\,y_i-y_{i-1})+1\bigr).\]

\begin{lema}[Sonda de rachas]\label{lem:sonda}
\begin{enumerate}[label=(\alph*),nosep]
\item $\Mrg(X,Y)$ tiene a lo sumo $\mathit{sw}$ cambios de lado (posiciones consecutivas
  de la salida que vienen una de $X$ y otra de $Y$).
\item La sonda hace a lo sumo $8\lceil\lg(b+1)\rceil$ comparaciones y no escribe.
\item Elegir una u otra mezcla según $\mathit{sw}$ no altera la poscondición del
  Lema~\ref{lem:sym}.
\end{enumerate}
\end{lema}
\begin{proof}
(a) Para $1\le i\le8$: todo $y$ de $Y$ antes de $y_i$ cumple $y\lt A[x_i]$ (los $y_i$ son
no decrecientes e $Y$ está ordenado), así que precede a $A[x_i]$ en $\Mrg(X,Y)$; todo $y$
desde $y_i$ cumple $y\nlt A[x_i]$ y, como en empates va primero $X$
(Def.~\ref{def:merge}), queda después de $A[x_i]$. Luego
$\Mrg(X,Y)=\cat_{i=1}^{9}\Mrg(\seg{A}{x_{i-1}}{x_i},\seg{A}{y_{i-1}}{y_i})$. Una mezcla
de $g_x$ y $g_y$ elementos tiene a lo sumo $2\min(g_x,g_y)$ cambios de lado, y entre dos
tramos consecutivos hay a lo sumo uno más.
(b) Son 8 búsquedas binarias sobre sufijos de $Y$.
(c) Las tres mezclas cumplen la misma poscondición: Lemas~\ref{lem:parity},
\ref{lem:mizq} y~\ref{lem:mizqg}.
\end{proof}

La regla es: galope si $8\,\mathit{sw}\le a+b$ (por (a), las rachas miden 8 o más en
promedio); si no, \Parity cuando $a+b\le\Bsz$, y el lazo clásico con saltos cuando no
entran juntos. Con menos de 256 elementos se usa \Parity sin sondear. El umbral y la
cantidad de muestras se eligieron midiendo en una sola máquina. Para los tipos que no se
copian byte a byte no hay sonda: se usa el galope, como en la versión 8. Por (b) y
porque \Natural hace a lo sumo $\lceil n/24\rceil-1$ mezclas, la sonda agrega
$O(n\log n)$ comparaciones en el peor caso y no cambia el Teorema~\ref{thm:cost}; con
mezclas balanceadas de tamaño $s\ge256$ hay $n/s$ por nivel y el total es $O(n)$.

\subsection{v8-r: sonda de orden global y regla de ruta}\label{sec:ruta}
\PocosGiros manda a \Natural toda entrada con menos de $n/8$ giros. Eso es un error
cuando los runs tienen claves entremezcladas (por ejemplo, tramos ordenados de 63 claves
al azar): \Natural hace $\lg(\text{runs})$ niveles de mezclas que alternan elemento por
elemento. En v8-r, si hay pocos giros se toman 256 muestras, una por tramo de largo
$\lfloor n/256\rfloor$ en una posición pseudoaleatoria del tramo, y se cuentan los
cambios de dirección entre muestras vecinas. Con menos de 64 cambios la entrada tiene
orden global y va a \Natural. Si no, va a \Natural solo cuando
$k\lceil\lg R\rceil\le\lceil\lg(n/\leaf)\rceil$, con $R=\lfloor g/2\rfloor+1$ la
estimación de runs a partir de los $g$ giros y $k$ el costo relativo de un nivel de
mezcla ($k=3$ o $2$ con buffer chico, según haya o no partición vectorial; $2$ o $5/4$
con buffer de al menos $n/2$). La sonda y la regla se aplican solo a tipos que se copian
byte a byte: con \emph{strings} mandar esas entradas al quicksort midió peor.

\begin{observacion}\label{obs:ruta}
La regla solo elige entre \Natural y \Quick. Los Teoremas~\ref{thm:natural}
y~\ref{thm:quick} valen para toda entrada que cumple sus precondiciones, sin importar
cómo se eligió el camino, así que el Teorema~\ref{thm:main} no cambia. La sonda hace
255 comparaciones y no escribe. Las constantes $k$ y el umbral de 64 salen de mediciones
en una sola máquina; una mala elección cuesta tiempo, nunca corrección. Sin el
desplazamiento pseudoaleatorio, runs de largo fijo producen muestras en posiciones
relativas que avanzan en progresión aritmética, y la entrada parece ordenada
(lo vi con runs de 1000 y $n=10^6$).
\end{observacion}

\subsection{v8-r: zonas desordenadas en MergeNatural}\label{sec:zonas}
Una entrada con orden global puede tener una zona al azar (el 10\% final, o un bloque
en el medio). Ahí \Natural encuentra runs de uno o dos elementos, los extiende a 24 con
inserción y los mezcla con mezclas mal predichas. En v8-r, cuando \Natural encuentra un
run de menos de 24 elementos mira hacia adelante: si ese run y los siguientes suman 8
o más runs seguidos de menos de 8 elementos, extiende la zona $[i,e)$ hasta el primer
run de 16 o más (o el final), la ordena con \Quick y la apila como un solo run. Si no,
hace lo de antes. Igual que la regla de ruta, se aplica solo a tipos que se copian byte
a byte.

\begin{lema}[Zonas]\label{lem:zonas}
Con esa modificación, el Teorema~\ref{thm:natural} sigue valiendo, y la pila de runs y
la cota de mezclas por elemento del Lema~\ref{lem:natcost} no cambian. El recorrido
hacia adelante agrega a lo sumo $3n$ comparaciones en total.
\end{lema}
\begin{proof}
El invariante de \Natural pide que cada entrada de la pila sea un intervalo ordenado y
coherente con los elementos de $\Omega$ restringido a él, y que los intervalos sean
adyacentes. La zona $[i,e)$ es un intervalo todavía no apilado; su contenido es
coherente porque solo se invirtieron runs estrictamente descendentes
(Lema~\ref{lem:yaord}). Por el Teorema~\ref{thm:quick}, \Quick deja $[i,e)$ ordenado y
coherente, que es lo que el invariante pide de un run. El buffer está libre entre dos
mezclas, así que \Quick puede usarlo. El Lema~\ref{lem:natcost} solo usa los largos de
los runs apilados, no cómo se obtuvieron, y la zona mide al menos 8 elementos. Costo
del recorrido: cada run mirado se recorre hasta 16 elementos como máximo, se miran a lo
sumo 8 runs antes de decidir que no hay zona (menos de $7\cdot 8+16=72$ comparaciones), y
después de decidirlo \Natural avanza al menos 24 posiciones; cuando hay zona, el
recorrido es lineal en la zona.
\end{proof}

El tiempo de la zona es el de \Quick sobre $e-i$ elementos; el Teorema~\ref{thm:cost}
no cambia porque \Quick está dentro de la misma cota.

\subsection{v8-r: recorridos vectoriales}
Para enteros con el orden natural, \Ordenado y \PocosGiros comparan $w$ pares adyacentes
por instrucción. Calculan exactamente los mismos predicados
($A[i]\lt A[i-1]$ para cada $i$) en el mismo orden de posiciones, así que los
Lemas~\ref{lem:yaord} y la Observación~\ref{obs:heur} no cambian. Las lecturas van de
$A[i-1]$ a $A[i+w-1]$ con $i+w\le r$: no salen del rango.

\subsection{Verificación de v8-r}\label{sec:v8r-pruebas}
Pruebas diferenciales contra \texttt{std::stable\_sort}, con el índice original como
etiqueta: enteros de 32 y 64 bits con y sin signo, pares, registros de 16, 40 y 128
bytes y \emph{strings}; los tres modos de memoria; de 1 a 4 hilos; orden por clave;
15 familias de entradas, incluidas zonas desordenadas y runs de claves al azar. Sobre el
archivo final: sin fallas en 19\,506 casos con AVX-512 y en 7\,410 con AVX2, sin SIMD y
con clang. AddressSanitizer, UBSan y ThreadSanitizer no informan nada. Se inyectaron
fallas de reserva en cada punto. Un contador en una versión anterior del mismo día
confirmó que las pruebas ejecutan las zonas (207 veces en una corrida), la regla de ruta
(480) y la sonda de orden (1245). Limitación conocida: si el
comparador lanza una excepción se pueden perder elementos; no hay errores de memoria.

%======================================================================
\section{Planificación fuera de orden (experimental)}\label{sec:ooo}

Esta sección estudia ejecutar los rangos pendientes de \Quick como tareas, en cualquier
orden y en varios hilos, sin las rondas de mezcla de \Paralelo. El prototipo
(\texttt{sotsort\_ooo.hpp}) está aislado de la biblioteca. Los resultados (a) a (e) del
Teorema~\ref{thm:ooo} están demostrados a mano al nivel de este documento; la
implementación del planificador solo está verificada por pruebas.

\subsection{Modelo}
Una \emph{tarea} es un rango pendiente $(l,r,\beta,\mathit{lb})$ de \Quick. Ejecutarla
es correr el lazo interior del Algoritmo~\ref{alg:quick} sobre ese rango, con una
diferencia: cuando el algoritmo secuencial apilaría la mitad grande de una partición,
la tarea puede \emph{cederla}: la pone en una cola compartida y sigue con la mitad
chica. La cola admite a lo sumo $64P$ tareas; con la cola llena la mitad grande se
apila localmente, como en el algoritmo secuencial. $P$ trabajadores toman tareas de la
cola hasta que está vacía y ninguno trabaja. Ninguna tarea espera a otra.

La \emph{huella} de una tarea es el conjunto de celdas que lee o escribe. Dos tareas
\emph{conmutan} si ejecutarlas en cualquier orden, o intercaladas, deja el mismo estado.

\begin{lema}[Conmutación]\label{lem:conm}
Si las huellas de $T_1$ y $T_2$ son disjuntas, $T_1$ y $T_2$ conmutan, y bajo el
Supuesto~\ref{sup:hilos} pueden ejecutarse a la vez sin carreras.
\end{lema}
\begin{proof}
Es el caso más simple de las condiciones de Bernstein: ninguna escribe lo que la otra
lee o escribe, así que cada una observa el mismo estado inicial en cualquier
intercalado y sus efectos se componen.
\end{proof}

\begin{lema}[Buffers disjuntos]\label{lem:bufdisj}
En modo completo, si mientras \Quick procesa $[l,r)$ se usa $\buf=\mathit{aux}+l$ y
$\Bsz=r-l$, toda escritura en el buffer cae en $\mathit{aux}[l..r)$. En los otros modos,
con un buffer por trabajador, toda escritura en el buffer cae en el buffer del
trabajador que ejecuta la tarea.
\end{lema}
\begin{proof}
Segundo caso: inmediato. Primer caso, por lo que se ejecuta con ese rango.
\Part: como $r-l\le\Bsz$ hay un solo bloque; \PartBloque escribe $\buf[i-m]$ con
$i<r-l$ y \PartVec escribe $\buf[i-m..i-m+w)$ con $i+w\le r-l$
(Lema~\ref{lem:partvec}); no hay fusiones, así que \Rotar no se ejecuta. \Hoja: escribe
$\buf[0..c)$ y $\buf[0..t)$ con $c\le t\le r-l$. \Pivote: a lo sumo $\Bsz$ muestras.
\Respaldo: \SymMerge sobre un subrango $[l',r')$ escribe desde $\buf[0]$ a lo sumo
$r'-l'\le r-l$ celdas.
\end{proof}

\begin{teorema}[Planificación fuera de orden]\label{thm:ooo}
Sea $A_0$ coherente con $n\ge2$ y supongamos los
Supuestos~\ref{sup:orden}--\ref{sup:hilos}. Para cualquier orden de ejecución de las
tareas y cualquier $P\ge1$:
\begin{enumerate}[label=(\alph*),nosep]
\item en todo momento, los rangos de la cola, los de las pilas locales y los que están
  en curso son disjuntos dos a dos y cada uno es pendiente en el sentido de la
  Def.~\ref{def:invq}; el resto de $[0,n)$ está terminado;
\item dos tareas en curso tienen huellas disjuntas;
\item las dependencias entre tareas forman un árbol;
\item la ejecución termina, sin bloqueo mutuo ni inanición;
\item al terminar, $A=\Omega$: la salida no depende del orden de ejecución ni de $P$, y
  es la de \SotSort.
\end{enumerate}
\end{teorema}
\begin{proof}
(a) El invariante de \Quick (Def.~\ref{def:invq}) habla de un conjunto de intervalos
pendientes; la demostración del Teorema~\ref{thm:quick} nunca usa que estén en una
pila ni el orden en que se procesan. Cada paso de una tarea es un paso del lazo
secuencial sobre su propio rango y preserva el invariante por la misma prueba. Ceder
una mitad en lugar de apilarla no cambia el conjunto.
(b) Por (a) los rangos son disjuntos; por el Lema~\ref{lem:bufdisj} los buffers
también. El estado del planificador se lee y escribe con un cerrojo tomado, y la cota
inferior viaja por copia (Supuesto~\ref{sup:ident}).
(c) Una tarea depende solo de la partición que la creó, y su rango está contenido
estrictamente en el de esa partición.
(d) Ninguna tarea espera a otra, así que no hay ciclo de espera. Cada tarea termina
por el Lema~\ref{lem:prof}. Las tareas son finitas: los rangos cedidos son disjuntos o
anidados y cada uno mide al menos el grano. La cola se atiende en orden de llegada.
Un trabajador termina solo cuando la cola está vacía y ningún otro trabaja, y en ese
estado no puede aparecer otra tarea.
(e) Al terminar no quedan pendientes, así que por (a) todo $[0,n)$ está terminado:
$A=\Omega$. $\Omega$ es única (Def.~\ref{def:omega}).
\end{proof}

\begin{observacion}
(e) es unicidad de la permutación estable, no una propiedad del planificador: cualquier
ordenamiento estable correcto da la misma salida. Lo que sí depende del diseño es que
las comparaciones realizadas no dependan del orden de ejecución: la semilla de
\Aleatorio de cada tarea cedida se deriva de $(\sigma,l,r)$, y la decisión de ceder
depende solo de los tamaños. Si la cola se llena, la decisión pasa a depender del
orden de ejecución; en las mediciones no se llenó nunca.
\end{observacion}

\subsection{Partición paralela de la raíz}
La primera partición es secuencial y fija el camino crítico. En modo completo se
reemplaza por tres fases. Sean $d_0=0<d_1<\dots<d_P=n$ y $P(x)\equiv x\lt p$.
\textbf{Fase 1} (en paralelo, tramo $i$): los que cumplen $P$ quedan en
$A[d_i..d_i+m_i)$ y el resto en $\mathit{aux}[d_i..d_i+f_i)$, en su orden (es
\PartBloque sin la copia final). \textbf{Fase 2} (secuencial, $i=1,\dots,P-1$):
$A[d_i..d_i+m_i)$ se mueve a $A[M_i..)$ con $M_i=\sum_{j<i}m_j$. \textbf{Fase 3} (en
paralelo): $\mathit{aux}[d_i..d_i+f_i)$ se copia a $A[M+F_i..)$, con $M=\sum_j m_j$ y
$F_i=\sum_{j<i}f_j$.

\begin{lema}[Raíz paralela]\label{lem:raiz}
Al terminar, $A=\filt(P,A_0)\cat\filt(\neg P,A_0)$, y las fases 1 y 3 no tienen
carreras.
\end{lema}
\begin{proof}
Fase 1: el tramo $i$ usa $A[d_i..d_{i+1})$ y $\mathit{aux}[d_i..d_{i+1})$, disjuntos
entre tramos. Fase 2: $M_i\le d_i$, así que el destino empieza a la izquierda del origen
y un movimiento hacia adelante es seguro; el destino termina en
$M_i+m_i\le d_i+m_i\le d_{i+1}$, antes del origen del tramo siguiente. Queda
$A[0..M)=\filt(P,A_0)$ porque cada tramo conserva su orden y los tramos se concatenan en
orden. Fase 3: los destinos $[M+F_i,M+F_i+f_i)$ son disjuntos, están en $A[M..n)$, y los
orígenes están en $\mathit{aux}$; $A[M..n)$ ya no tiene nada pendiente.
\end{proof}

\subsection{Trabajo y camino crítico}
El trabajo $W(n)$ es el de \Quick más $O(1)$ por tarea. Con la partición secuencial, el
camino crítico es $S(n)=\Theta(n)$ esperado (la partición de un rango de largo $\lambda$
cuesta $\Theta(\lambda)$ y, siguiendo la mitad grande, las $\lambda$ decrecen
geométricamente), así que el paralelismo $W/S$ es $\Theta(\log n)$. Con la raíz paralela
el primer término baja a $n/P+(M-m_0)+\max_i f_i$. Dentro de \Part por bloques, con
rotaciones secuenciales, la fusión de la raíz mueve $\Theta(\lambda)$: el paralelismo
interno es $O(\lg(\lambda/\Bsz))$, y por eso no se construyó una variante con tareas por
bloque.

\subsection{Qué no está demostrado}
\begin{itemize}[nosep]
\item Que la implementación del planificador (cola, cerrojo, variable de condición)
  cumple el modelo. Evidencia: 466 casos contra \texttt{std::stable\_sort} con 2 a 5
  hilos y grano de 64 a 32\,768, y 60 corridas con un comparador que lanza, todo sin
  avisos de ThreadSanitizer ni de AddressSanitizer.
\item Cualquier afirmación de escalabilidad: solo se midió con 2 núcleos.
\item El modelo de memoria de C++ en detalle; el Supuesto~\ref{sup:hilos} lo resume.
\end{itemize}


%======================================================================
\section{Limitaciones y trabajo futuro}
\begin{itemize}[nosep]
\item Formalizar mecánicamente estas pruebas (Why3 o Dafny para las triplas de Hoare,
  Coq o Isabelle para los lemas de órdenes y de powersort).
\item Probar la equivalencia con el código C++, o generar el código desde la
  especificación verificada.
\item Las afirmaciones de rendimiento contra blitsort son experimentales: una sola
  máquina, ruido de $\pm10\%$ y solo \texttt{uint64}.
\item Créditos: la triple reversión fusionada sigue la \emph{trinity rotation} de
  blitsort; el \emph{parity merge} de igual largo y el \emph{cross merge} vienen de
  quadsort y blitsort (scandum, dominio público); el criterio de powersort es de
  Munro y Wild, con la formulación del \texttt{powerloop} de CPython. No se hizo una
  revisión de literatura exhaustiva.
\end{itemize}

\end{document}
